\documentclass[11pt,a4paper]{article}

\usepackage[margin=1in]{geometry}
\usepackage{amsmath}
\usepackage{amssymb}
\usepackage{amsthm}
\usepackage{booktabs}
\usepackage[numbers]{natbib}
\usepackage[hidelinks,unicode]{hyperref}

% plainnat emits \doi{...} for article and inproceedings entries and provides a plain-text fallback. Define it first
% so every DOI becomes a resolvable link, and typeset the identifier URL-style so raw underscores in DOI values pass
% through unescaped; techreport and misc entries never reach \doi from the style, so their notes call it directly.
\newcommand{\doi}[1]{\href{https://doi.org/#1}{\nolinkurl{doi:#1}}}

\hypersetup{
  pdftitle={Certified Split Windows for Parallel Lexing: Recovering Boundaries Where No Byte Certifies},
  pdfauthor={Nicklas Nidhögg},
  pdfsubject={Parallel lexical analysis; deterministic finite automata; maximal munch tokenization},
  pdfkeywords={parallel lexing, tokenization, deterministic finite automata, maximal munch, split windows}
}

\theoremstyle{definition}
\newtheorem{definition}{Definition}
\newtheorem{remark}{Remark}

\theoremstyle{plain}
\newtheorem{lemma}{Lemma}
\newtheorem{theorem}{Theorem}
\newtheorem{corollary}{Corollary}
\newtheorem{proposition}{Proposition}
\newtheorem{invariant}{Invariant}

\newcommand{\code}[1]{\texttt{#1}}
\newcommand{\before}{\mathsf{before}}

% The default itemize bullet pulls in a Type 3 bitmap font, which some submission systems reject. The math bullet
% comes from the Type 1 symbol font that is embedded already.
\renewcommand{\labelitemi}{\(\bullet\)}

\title{Certified Split Windows for Parallel Lexing:\\Recovering Boundaries Where No Byte Certifies}
\author{Nicklas Nidh\"ogg\\[0.4em] \normalsize Independent Researcher\\ \normalsize
\texttt{nicklas.nidhogg@gmail.com}\\ \normalsize ORCID: 0009-0006-6161-6150}
\date{August 2026}

\begin{document}

\maketitle

\begin{abstract} A certified split point lets a parallel lexer cut unlexed input at a single byte with the serial token
stream provably preserved. Yet several conventional token sets in the predecessor's controlled study certify no byte
once string, comment, or whitespace-run forms are included (arXiv:2608.03473). We generalize from a byte to a bounded
window: a byte string after which the current token's start is known, regardless of surrounding context. We certify the
directly usable form of that recovery: the token covering the window's final byte begins at the reported origin. The
certificate is conditional on occurrence and may be vacuous. Every applicability figure counts only windows with an
asserted completely tokenizable occurrence witness. We give a conservative model of a maximal-munch scanner's possible
histories across a window and prove it sound. We decide reachability in that model exactly by exhausting a finite
quotient of its reachable configurations, so every answer of the unbudgeted procedure is a certified window with its
origin or proof that the model admits none. Within the stated flat, completely-tokenizable scope, model-positive answers
are semantic certificates. Negatives are relative to the model, which deliberately refuses some windows a greedy scanner
would allow. A token set with a token matching the empty string is decided through its positive-width equivalent, the
same automaton entered at a non-accepting start state, which changes no scan. Of 400 random token sets, 322 of the 337
certifying no byte gain a witnessed window, with zero inconclusive searches. Every exact-empty row of the predecessor's
study gains a witnessed window of two to four bytes. The analysis runs once after automaton construction, using only the
compiled tables and no input. A boundary profile decides each gap of a window, forced, crossed or open; a window-local
cut is sound exactly at forced gaps. Existence of a model certificate, plain or witnessed, is PSPACE-complete, by a
reduction from DFA intersection non-emptiness that fixes the minimum length at the source's plus a constant, and a
shortest window can need $2^{\Omega(\sqrt n)}$ bytes over $n$ states. Two smaller searches decide it, a quotient of
$3^n$ keys and an unambiguous automaton of $(n + 2)\, 2^{n-1}$ states. The second is the search that release 2.2 of
munch ships. The cloud's hypotheses
are exactly the histories of arbitrary factorizations of words of the token language, which explains the vacuous
window, makes the model exact on a polynomially decidable class of token sets strictly beyond the prefix codes, and
makes it decidable whether the model misses any witnessed certificate at all.
\end{abstract}

\section{Introduction}
\label{sec:introduction}

The predecessor of this paper derives, from a compiled token set, the complete set of bytes at which unlexed input can
be cut with the serial token stream preserved. The predecessor proves the condition necessary as well as
sufficient~\cite{nidhogg2026splitpoints}. Its sharpest limitation is its own applicability table. Several of the studied
conventional token sets certify \emph{no} byte, because a single string form, comment form, or whitespace run gives some
non-initial live state a transition on every candidate. A fresh 400-grammar random sweep generated for this study
reproduces the pattern, 337 of its 400 sets certifying none. This paper is about the object that refusal leaves
standing: a byte \emph{string} after which the start of the token covering the window's final byte is pinned, regardless
of surrounding context.

Concretely, take a token set compiled to a DFA and scanned by maximal munch. We ask whether there is a window $W = w_0
\cdots w_{k-1}$ and an offset $o$ with $0 \leq o < k$ such that in \emph{every} completely tokenizable input containing
$W$, the token covering the occurrence's final byte begins exactly $o$ bytes into it. A worker that finds $W$ in a
completely tokenizable input may then begin scanning at the recovered boundary. Exactly as at a certified byte, which is
the $k = 1$ case, the worker needs no speculation, no state enumeration, and no DFA-state-recovery pass over the input.

Our contributions:

\begin{itemize}
\item
A conservative model of the scanner's possible token-prefix histories across a window, replacing a refuted
natural predecessor, and a soundness proof by a representation
invariant
(Section~\ref{sec:model}, Section~\ref{sec:soundness}).
\item
A finite quotient of the model's reachable configurations, exact for them by a stated single-occupancy invariant,
turning breadth-first search over windows into a terminating decision procedure for the model
(Section~\ref{sec:quotient}).
\item
Two specializations (Section~\ref{sec:specialization}): at length one the model certifies exactly the bytes the
published predicate reports, so the multi-byte construction is that predicate's conservative continuation. Over a prefix
code the certified windows that occur are exactly the synchronizing splits of code theory whose right half sits inside
one codeword, which places the classical case as the special case it is.
\item
Two smaller exact searches for the model (Section~\ref{sec:smaller}): a quotient of at most $3^n$ keys over $n$ live
states, and an unambiguous automaton of at most $(n + 2)\, 2^{n-1}$ states that follows one chosen origin. The second
returns a shortest certified window, bounds its length, and is the search release 2.2 of munch ships as
\code{shortest\_split\_window()}.
\item
The complexity of the model (Section~\ref{sec:complexity}): existence of a model certificate, plain or witnessed, is
PSPACE-complete, by a reduction from DFA intersection non-emptiness with one token kind and five symbols that fixes the
minimum window length at the source's shortest common word plus $m + 3$. A family of such token sets has shortest
windows of length $2^{\Omega(\sqrt n)}$, with exact minima 4, 6, 11, 18, 67, 68, 429, 850 and 2{,}531 for its first nine
members, searched by a figure program.
\item
What the cloud represents (Section~\ref{sec:factorization}): its hypotheses are exactly those of arbitrary
factorizations of words of the token language into tokens, greedy or not. The vacuous window follows, token kinds
may be erased before any search, and the model is exact wherever every factorization is the scan's, a class decided
in polynomial time that holds the prefix codes strictly. A shortest witnessed model certificate is found exactly with
its occurrence, by a product with the companion's verifier, and whether the model misses any witnessed semantic
certificate at all is decidable, with a shortest missed window.
\item
An evaluation over all six exact-empty rows of the prior study's applicability table, one new cumulative variant, and
400 random token sets generated for this paper (Section~\ref{sec:evaluation}). The evaluation reports the witnessed
rescue rate where no byte certifies and the window lengths that suffice for the studied C-like and JSON tokenizations.
Rewind-stress checks of 1{,}079{,}392 generated executions that scanned through the window and contained at least one
rewind, 418{,}466 of them completely tokenizable, found zero disagreements against the shipped scanner.
\item
A boundary profile of a window, the verdict forced, crossed or open at each of its gaps over every occurrence, with the
certificate as one of its patterns (Section~\ref{sec:profile}). A window-local cut rule that reconciles nothing is sound
exactly at forced gaps. Over a token set with a longest token, a window containing a forced gap forces every boundary
the scan from that gap commits with its lookahead inside the window. No single width decides every gap of the JSON
tokenization, since no bounded window sees whether an unmatched quote lies to its left. Release 2.1.0 of munch decides
the profile exactly within its search cap, and an exploratory campaign over JSON, C-like and byte-pair token sets
measures what it adds to the certificate.
\end{itemize}

The probe decides certificates and model-search results offline from the compiled tables after automaton construction.
The occurrence witnesses and scanner checks are generated executions against the shipped scanner. Every empirical
aggregate and table entry of the evaluation (Section~\ref{sec:evaluation}) is printed and asserted by the probe in the
munch repository~\cite{munch}, so a drifted number fails the test suite. The named rows are asserted at release v1.3.3,
and the random sweep at the later commit Section~\ref{sec:evaluation} names. The boundary-profile measurements of
Section~\ref{sec:profile} are an exploratory campaign the probe does not assert.

\section{Preliminaries}
\label{sec:preliminaries}

We inherit the scanner model of~\cite{nidhogg2026splitpoints}. A token set compiles to a DFA $A$ with initial state
$q_0$. Scanning is by maximal munch: from the current position the scanner runs $A$ as far as a transition exists and
emits the token of the last accepting configuration passed. The scanner resumes immediately after that token and
restarts in $q_0$. An input is \emph{completely tokenizable} when this process consumes it exactly. The setting is
\emph{flat}: one fixed token set compiled to one DFA scanned from one fixed $q_0$, with no lexical modes, no mode stack,
and no semantic scanner state. $A^+$ denotes the live subautomaton: states both reachable from $q_0$ and co-accessible
to acceptance, with only transitions between live states retained. A token may match the empty string, and the scan
never emits it. The next lemma is how every statement below still assumes a start state that does not accept.

\begin{lemma}[Positive-width equivalent]
\label{lem:unroll}
Let $A$ be the DFA of a token set whose start state $q_0$ accepts, and let $A'$ be $A$ with a fresh start state
$q_0'$ that carries $q_0$'s outgoing transitions and does not accept, $q_0$ itself retained. Then the maximal-munch
scans of $A$ and $A'$ agree on every input, $A'$ accepts exactly the nonempty words $A$ accepts, and no transition of
$A'$ enters $q_0'$.
\end{lemma}

\begin{proof}
Under the scan's convention only an accepting configuration reached after consuming a byte determines an emitted token,
so every emitted token has positive width and a scan records an acceptance only after consuming a byte. From $q_0$ and
from $q_0'$ the first byte leads to the same state, since $q_0'$ carries $q_0$'s transitions. From there the two scans
traverse the same states and record the same accepting positions. The emitted token and the restart offset therefore
agree, and both scans restart in their own start state. The accepted words of $A'$ are the words of length at least one
accepted by $A$, the empty word being lost because $q_0'$ does not inherit $q_0$'s acceptance. No transition of $A$
targets $q_0'$, which is fresh, and the copied transitions target $q_0$'s successors, so nothing enters $q_0'$.
\end{proof}

Throughout, a token set whose start state accepts is read through $A'$. The artifact compiles every such token set this
way before deciding anything about it, and a set whose start does not accept passes through unchanged. The unrolled
automaton has one state more and the same scan, so nothing below loses generality by assuming that $q_0$ does not
accept. From here on $A$ and $q_0$ name the automaton so read and its start, $q_0'$ for a nullable set. For a nullable
set the fresh start is never re-entered. For any other set, the re-entrancy of $q_0$ remains the separate condition
Section~\ref{sec:specialization} treats.

\begin{definition}[Certified split window]
\label{def:window}
Let $W = w_0 \cdots w_{k-1}$ with $k \geq 1$ and let $o \in \{0, \ldots, k-1\}$. The pair $(W, o)$ is a
\emph{certified split window}
for a token set when, for every completely tokenizable input $x$ containing $W$ at offset $t$, the token of the
maximal-munch tokenization of $x$ that contains the byte at offset $t + k - 1$ begins at $t + o$.
\end{definition}

\begin{definition}[Witnessed window]
\label{def:witnessed}
A certified split window $(W, o)$ is \emph{witnessed} when some completely tokenizable input contains $W$.
\end{definition}

Definition~\ref{def:window} quantifies over the inputs containing $W$ and is therefore vacuously true when no completely
tokenizable input contains $W$. The vacuity is not hypothetical. Over $\{\code{0}, \code{00}, \code{01}\}$ the model
below certifies the window \code{1001} at origin 2, yet no completely tokenizable input contains \code{1001}. The reason
is that every \code{1} is the tail of a greedily chosen \code{01}, so a token boundary follows the window's first byte.
Maximal munch must then consume \code{00}, and the final \code{1} sits at a boundary no token starts. The preceding
argument proves non-occurrence. Section~\ref{sec:factorization} explains the certificate: \code{1001} occurs in a
factorization of \code{01001} into tokens that the scan does not take. The artifact asserts the model prediction and the
empty result of its bounded targeted witness search. Every applicability figure in this paper counts only witnessed
certificates. The named rows pin their witness inputs in the artifact. For each of the 322 witnessed random grammars the
probe constructs and verifies a concrete completely tokenizable input, with the 322 aggregate asserted. A certificate
the bounded witness search cannot witness is reported as unresolved, never as a rescue.

For $k = 1$ Definition~\ref{def:window} is the boundary guarantee of a certified split point, since the token containing
the single byte begins at it. Three guarantees should be kept apart, and this paper certifies the strongest. Model
unanimity implies that the covering token's origin is fixed, which implies that some fixed safe boundary exists inside
the window, and neither converse holds. Over $\{\code{a}, \code{ab}, \code{b}\}$ at \code{ab} the covering origin is
fixed yet the model refuses (Section~\ref{sec:strictness}). Over $\{\code{a}, \code{abx}, \code{b}, \code{x}\}$ at
\code{ab} the byte \code{a} always begins a token, so offset $0$ is a fixed safe boundary. Yet the covering token's
origin there is not fixed, since the final \code{b} belongs to \code{b} in the input \code{ab} and to \code{abx} in
\code{abx}. The weakest guarantee is already usable, since a worker can cut at the known boundary and rescan the
window's suffix. Release 2.1.0 of munch decides that weakest guarantee exactly within its search cap
(Section~\ref{sec:profile}): its \code{boundary\_profile()} classifies every gap of an occurring window as \code{must},
\code{never} or \code{may}, a boundary at every occurrence, at none, or at some. The covering-token form is the
sufficient, deliberately stronger property this forward origin-recovery model certifies, and that form hands the worker
its resumption point directly. The model therefore has two distinct sources of false negatives: covering-origin
certification is stricter than locating some safe boundary, and the cloud conservatively represents extra segmentations
and may refuse even a true covering-origin certificate. Note also what the definition does not require: it says nothing
about the tokens overlapping the window's earlier bytes, and it does not require the boundary to be the only one inside
the window.

\section{The model}
\label{sec:model}

A worker cutting blind knows only that in the final segmentation, the window's first byte is consumed from \emph{some}
live token-prefix state, either inside a token that began at some unknown earlier offset or exactly at a boundary. The
acceptance-gated seed below carries the boundary case, including a window at the start of the input. The model tracks a
\emph{cloud} of hypotheses $(q, \omega)$: a live state $q$ paired with an origin $\omega$, either $\before$ for a token
that began before the window or an in-window offset. The initial cloud $C_0$ is every live state paired with $\before$.

Reading window byte $w_j$ maps $C_j$ to $C_{j+1}$ by two rules:

\begin{itemize}
\item \textbf{Direct step.} A pair $(q, \omega)$ whose state consumes $w_j$ into a live state moves there with its
origin unchanged. A pair whose state cannot consume $w_j$ into a live state is an impossible history and is dropped,
never restarted.
\item \textbf{Acceptance-gated seed.} If some pair in $C_j$ is accepting, one fresh pair $(\delta(q_0, w_j),\, j)$ is
seeded, provided that target is live. The gate is sound because tokens end only where the automaton accepts, and a token
can begin at offset $j$ only where the previous token could have ended or at the input start. The input start is the
case the initially open gate represents.
\end{itemize}

One rename accompanies the direct step. Where $q_0$ is not re-entrant in $A^+$, a pair stepping \emph{from} $q_0$ is
beginning a token, so its origin becomes the current offset. Where a non-empty live path returns to $q_0$, the rename is
disabled, since reaching $q_0$ then no longer identifies a boundary. The length-one certificate carries the same
re-entrancy condition~\cite{nidhogg2026splitpoints}.

Clouds are \emph{sets} of pairs. The set semantics is load-bearing below, where two rules inserting the same pair yield
one element. The window is \emph{certified at origin $o$} when $C_k \neq \varnothing$ and there is an $o \in \{0,
\ldots, k-1\}$ with $\omega = o$ for every $(q, \omega) \in C_k$. An empty cloud certifies nothing, matching the
implementation, which refuses it. The empty cloud is absorbing: every step of it is empty. Non-emptiness guarantees
nothing in the other direction: a certified window may be vacuous under Definition~\ref{def:window}, the separation
Definition~\ref{def:witnessed} exists to make. Agreement on the state alone is not enough: learning that the scan is
inside a string literal is knowledge, but not a boundary. Throughout, $q_0$ is \emph{re-entrant} when some live
transition of $A^+$ targets it. Non-re-entrancy means no live edge enters $q_0$ at all.

\paragraph{The discarded predecessor.} A natural variant restarts a trajectory that cannot consume the byte, treating
the failure point as a token boundary, in place of the acceptance-gated seed. That variant is not merely unproved but
refuted. Over the token set $\{\code{a}, \code{abc}, \code{bx}, \code{x}\}$ and window \code{abx}, the variant's cloud
after \code{ab} is the single pair carrying origin 0 toward \code{abc}. The byte \code{x} kills that trajectory, and the
restart replaces it with a token beginning at offset 2. Nothing else survives, and the variant certifies $(\code{abx},
2)$. Yet the input \code{abx} itself is completely tokenizable as \code{a} followed by \code{bx}, so the token covering
the final byte begins at offset 1: the certificate is false at a witnessed occurrence. A failing trajectory is an
impossible history, not a boundary, and restarting it manufactures support for origins no execution justifies. The
repaired model certifies \code{abx} at origin 1, where the scanner does cut, and the case is carried as an asserted row
of the artifact.

\section{Soundness}
\label{sec:soundness}

Fix a completely tokenizable input $x$ containing $W$ at offset $t$. For $j \in 1..k$ let $\sigma_j$ be the start of the
token that the final maximal-munch segmentation of $x$ assigns to the byte at $t+j-1$. For the same $j$, let $\rho_j =
\delta^*(q_0,\, x[\sigma_j \,..\, t+j))$ be that token's prefix state after consuming through byte $t+j-1$, and let
$\omega_j = \sigma_j - t$, or $\before$ when $\sigma_j < t$.

\begin{lemma}[Representation]
\label{lem:representation}
For every $j \in 1..k$, the pair $(\rho_j, \omega_j)$ is in $C_j$.
\end{lemma}

\begin{proof}
Every $\rho_j$ is live: reachable through the actual token prefix, and co-accessible because its token ends at
some $e \geq t+j$ with $\delta^*(\rho_j, x[t+j \,..\, e)) = \delta^*(q_0, x[\sigma_j \,..\, e))$, and the right
side is accepting.

\emph{Base, $j = 1$.} If $\sigma_1 < t$, the state $p = \delta^*(q_0,\, x[\sigma_1 \,..\, t))$ is live and lies in $C_0$
paired with $\before$. The direct step carries $p$ to $(\rho_1, \before)$. The rename cannot interfere: $p$ has consumed
at least one byte, so $p = q_0$ only if a non-empty live path returns to $q_0$, which disables the rename. If $\sigma_1
= t$, then $\rho_1 = \delta(q_0, w_0)$ and the seed fires, because $C_0$, which is all of the live states, contains a
live accepting state. The fixed input is completely tokenizable and non-empty, so its first token traces an accepting
path from $q_0$. That path makes $q_0$ live with an accepting state reachable from it, and a reachable accepting state
is trivially co-accessible.

\emph{Step, $j$ to $j+1$.} If the byte at $t+j$ continues its token, $\sigma_{j+1} = \sigma_j$, and the direct step
carries $(\rho_j, \omega_j)$ to $(\delta(\rho_j, w_j), \omega_j)$. In this case the rename does not overwrite the
origin, by the same re-entrancy argument as in the base. If the byte at $t+j$ begins a token, the previous token ended
at $t+j$, so $\delta^*(q_0,\, x[\sigma_j \,..\, t+j))$ accepts. That accepting state is $\rho_j$, in $C_j$ by
hypothesis, so the seed fires and emits exactly $(\delta(q_0, w_j),\, j) = (\rho_{j+1}, \omega_{j+1})$.
\end{proof}

The lemma is containment, not equality: the cloud may carry hypotheses no execution realizes. The surplus is harmless in
one direction and load-bearing in the other. At a realized occurrence in a completely tokenizable input, surplus
hypotheses can only preserve the actual origin's unanimity or destroy unanimity. They cannot manufacture unanimity at a
false origin.

\begin{theorem}[Soundness]
\label{thm:soundness}
If $C_k \neq \varnothing$ and every pair of $C_k$ carries the same origin $o \neq \before$, then $(W, o)$ is a
certified split window.
\end{theorem}

\begin{proof}
By Lemma~\ref{lem:representation}, $(\rho_k, \omega_k) \in C_k$, so $\omega_k = o$, so $\sigma_k = t + o$, a token
start of the final segmentation. The input was an arbitrary completely tokenizable one containing $W$ at $t$.
\end{proof}

\paragraph{Backup never appears.} Maximal-munch lookahead that a later rewind discards occupies states the lemma says
nothing about. The invariant tracks where the \emph{final} segmentation's tokens begin, not where the read head wanders.
A boundary that a rewind later exposes was seeded by the model step reading the first byte after the accepting position
justifying it. That seeding argument holds when lookahead crosses the window's left edge, when several accepting
positions are passed, and across chains of consecutive rewinds. Each of these cases is an instance of the step case.

\section{The finite quotient and the decision procedure}
\label{sec:quotient}

Acceptance gating alone does not terminate: over $\code{a}^+$ the automaton accepts after every byte and origins
accumulate without bound. The search therefore deduplicates on a quotient of the cloud: the set $B$ of states
carrying $\before$, and for each state the number of distinct in-window origins it carries, saturated at two.

\begin{invariant}[Single occupancy]
\label{inv:single}
In every reachable cloud, each surviving in-window origin occupies at most one state.
\end{invariant}

An origin enters the cloud at most once per rule application, into a single state, and the deterministic step moves it
to at most one successor, so it continues in one state or dies. At offset $j$ the acceptance seed and the
non-re-entrant-$q_0$ rename can both propose the fresh origin $j$. Both target exactly $\delta(q_0, w_j)$, and set
semantics coalesces the identical pair, so single occupancy survives the collision. Origins may die, which is why the
invariant says at most one rather than exactly one. The pre-window origin is deliberately different: it may occupy many
states and is held apart as the Boolean support set $B$. An arbitrary cloud could violate the invariant, placing one
origin in two states, and there the saturated counts could not decide unanimity. No such cloud is reachable, and the
restriction is load-bearing for everything below. The key of a cloud is the pair $\kappa(C) = (B,\, m)$ with $m(q) =
\min(2,\, \text{number of in-window origins at } q)$.

\begin{lemma}[Congruence, depth-aligned] \label{lem:congruence} Let $C_u$ and $C_v$ be clouds reached by the search
after words $u$ and $v$. If $\kappa(C_u) = \kappa(C_v)$, then for every byte $b$ the successors $\mathrm{step}(C_u, b,
|u|)$ and $\mathrm{step}(C_v, b, |v|)$ are either both empty or both non-empty, and when non-empty, \[
\kappa(\mathrm{step}(C_u, b, |u|)) = \kappa(\mathrm{step}(C_v, b, |v|)), \] with the certification verdict agreeing on
both sides. For successor-key equality alone, it is enough that the inserted offsets be \emph{fresh}, exceeding every
in-window origin of their respective clouds; the certification verdict additionally needs legal depths, and the search
supplies both: it steps with the word lengths $|u|$ and $|v|$, and every origin in a length-$\ell$ cloud lies below
$\ell$. Freshness is not decorative: over $\{\code{a}^+, \code{b}\}$ the cloud after \code{a} carries origin $0$, and
stepping on \code{a} with offset $0$ merges the seed into that origin while offset $1$ creates a second one, so the
same key would have two different successors under arbitrary offsets. \end{lemma}

\begin{proof}
The acceptance gate reads only which occupied states accept, and the occupied states are $B \cup \{q : m(q) > 0\}$, a
function of the key. The new $\before$-support is the live deterministic image of $B$, excluding the contribution from
$q_0$ when $q_0$ is non-re-entrant. The excluded contribution joins the fresh-origin rename handled below. Over the
single-token set $\{\code{a}\}$, stepping $C_0$ on \code{a} renames $q_0$'s $\before$ origin to $0$, so $B' =
\varnothing$, where the unexcluded image would wrongly retain $\delta(q_0, \code{a})$. Consider the in-window counts at
a successor state $q'$. The direct step carries each origin from its unique state (Invariant~\ref{inv:single}), so
origins arriving at $q'$ from distinct predecessors are distinct. The exact update at each target sums the old-origin
contributions arriving there. The update adds one fresh origin if the seed or the non-re-entrant-$q_0$ rename fires
toward the target, combining those two by logical or, since they denote the same fresh origin and the cloud is a set.
The update finally saturates at two. When $q_0$ is non-re-entrant, no live edge enters $q_0$, so after the first step no
reachable cloud contains $q_0$ with an in-window origin. For a non-re-entrant start state the rename decision is
therefore readable from $q_0 \in B$. Whether the rename fires is determined by $q_0 \in B$, which is key-readable, and
by re-entrancy, which is fixed automaton data rather than part of the key. Whether the seed fires is determined by the
key's accepting support. The freshness hypothesis is exactly that the fresh origin collides with no surviving origin.
Saturation commutes with this update in the only direction needed. A saturated state maps its surviving origins
together, so a successor's saturated sum computed from saturated counts equals the saturation of the exact sum. A
state's count can also drop to zero outright when its transition is missing, which the update handles as an empty
contribution. Fresh origins are equivariant under renaming of origin values, and no rule ever reads an origin's value,
only equality between origins, so keys computed on either side agree. Emptiness is key-readable. Certification of the
full window is the key predicate $B = \varnothing$ with total count exactly one, using Invariant~\ref{inv:single} to
identify one total count with one origin in one state.
\end{proof}

One tempting simplification is false and worth flagging: it is not true that a state holding two origins can never
contribute to unanimity later. Over $\{\code{a}^+, \code{b}\}$, after \code{aa} the \code{a}-state carries two in-window
origins, as well as $\before$. On \code{b} all of them die while their pre-step acceptance opens the gate for one fresh
\code{b}-origin, and the cloud is unanimous. What is true, and what the lemma uses, is that co-located origins either
follow the same live successor together or all die, and that origins which die leave no trace the key must distinguish.

\begin{theorem}[Decision procedure]
\label{thm:decision}
Breadth-first search over windows, deduplicated on keys, decides whether the model certifies any window for a
given token set, and returns the minimum certified length; it does not enumerate every certified word. Each state
contributes a factor of $2 \times 3$ to the key space, so at most $6^{|Q^+|}$ keys are ever retained; each key
expands into at most $256$ successors, each computable from the key in $O(|Q^+|)$ time, for a worst case of
$O(256 \cdot |Q^+| \cdot 6^{|Q^+|})$ time and $O(|Q^+| \cdot 6^{|Q^+|})$ space.
\end{theorem}

By Lemma~\ref{lem:congruence} the walk cannot miss the existence of a certifying continuation, nor change the minimum
certified length. After one matched step both successor clouds sit at fresh depths again, so the lemma inducts over
every common suffix. Breadth-first order retains a shallowest representative of each key. A certifying suffix from a
discarded deeper occurrence therefore yields an equal or shorter certificate from the retained representative.
Deduplication may still skip individual certified words. Exhaustion proves the model admits none at any length. What
exhaustion does not give is semantic non-existence, because the model itself is conservative
(Section~\ref{sec:strictness}). The displayed bounds describe an abstract key-to-key search. The probe realizes the
procedure with one concrete representative cloud and word per key, using the quotient for deduplication, so the probe's
cost additionally depends on representative size. The probe is budgeted, with three outcomes: certified, exhausted, and
inconclusive once the retained key count exceeds a fixed threshold of 200{,}000 keys. The evaluation below reports zero
inconclusive searches. The retained key counts the evaluation reports are one indicator of the search footprint rather
than a complete cost model. Section~\ref{sec:smaller} replaces the $6^{|Q^+|}$ bound by $3^{|Q^+|}$ and by
$(|Q^+| + 2)\, 2^{|Q^+| - 1}$, and shows the exponential cannot be removed.

\section{Smaller searches and the complexity of the model}
\label{sec:smaller}

Write $n = |Q^+|$ for the number of live states. The quotient of Section~\ref{sec:quotient} is exact and finite. Its
$6^n$ bound is what the evaluation's retained-key counts stand against. Two smaller presentations decide the same
question. The first coarsens the key to three statuses per state. The second guesses the one origin that survives,
and is an unambiguous automaton over the certified windows. The second is the search release 2.2 of munch ships. It
is also the route to the model's exact complexity. Existence of a model certificate is PSPACE-complete, and a shortest
window can be exponentially long in $n$.

\subsection{A three-status quotient}
\label{sec:three}

Fix a reachable cloud. A live state is \emph{empty} when no pair occupies it. It is \emph{single} when exactly one
pair occupies it and that pair carries an in-window origin. Otherwise it is \emph{blocked}: it carries $\before$, or
at least two in-window origins. The three-status key of a cloud is the pair $\kappa_3(C) = (S, B)$ of its single and
its blocked states. A state carrying $\before$ beside an in-window origin is blocked. So $\kappa_3$ identifies keys
that the six-status $\kappa$ keeps apart.

\begin{theorem}[Three-status quotient]
\label{thm:three}
Let $C_u$ and $C_v$ be clouds reached by the search after words $u$ and $v$ with $\kappa_3(C_u) = \kappa_3(C_v)$. For
every byte $b$ the successors $\mathrm{step}(C_u, b, |u|)$ and $\mathrm{step}(C_v, b, |v|)$ are both empty or both
non-empty. When non-empty, their three-status keys agree. A reachable cloud is certified exactly when $B =
\varnothing$ and $|S| = 1$. Breadth-first search over three-status keys therefore decides whether the model certifies
any window and returns the minimum certified length. It retains at most $3^n$ keys, in $O(256 \cdot n \cdot 3^n)$
time and $O(n \cdot 3^n)$ space.
\end{theorem}

\begin{proof}
The occupied states are $S \cup B$, so the key says whether the acceptance gate is open. When $q_0$ is not
re-entrant, no live edge enters it. Then $q_0$ is occupied only in $C_0$, where it carries $\before$ and is blocked.
So the key says whether the rename fires, through $q_0 \in B$. When it fires, the contribution of $q_0$ to the direct
step is set aside and counted as the fresh origin, as in Lemma~\ref{lem:congruence}. Now fix a successor state $q'$
and collect what arrives there. Suppose a blocked predecessor $p$ has $\delta(p, b) = q'$. It brings $\before$, or
two distinct in-window origins, and all of them move to $q'$ together. So $q'$ is blocked, whatever else arrives.
Otherwise every predecessor of $q'$ is single and brings one in-window origin. These origins are pairwise distinct
by Invariant~\ref{inv:single}, since one origin occupies one state. The fresh origin, when the seed or the rename
targets $q'$, is distinct from all of them by freshness. So $q'$ is blocked, single or empty as the number of
arriving origins is at least two, one or zero. Every quantity used is a function of $(S, B)$, of $b$ and of fixed
automaton data. Hence the successor keys agree. The successor is empty exactly when every state is empty, which the
key says too.

For the verdict, a non-empty cloud is certified when no pair carries $\before$ and all pairs share one in-window
origin. No $\before$ and no two origins at one state means $B = \varnothing$. Then every occupied state is single. By
single occupancy the number of distinct origins is then the number of occupied states, so one origin means $|S| = 1$.
The congruence inducts over every continuation as in Theorem~\ref{thm:decision}, with the same depth alignment. The
bounds follow from the three statuses per state.
\end{proof}

The coarsening loses nothing, because a blocked state can never again supply the one permitted origin. Its origins
move together or die together, which is the observation that closed Section~\ref{sec:quotient}. The six-status key
additionally separates, at a blocked state, whether it carries $\before$ and how many in-window origins up to two,
and neither distinction is ever read. The figure program
of Section~\ref{sec:evaluation} runs this search beside the gate's. It finds the same shortest length on every named
row, and Table~\ref{tab:windows} gives the retained keys of both.

\subsection{The chosen-origin search}
\label{sec:chosen}

A certified window has one surviving origin. Guess that origin at its birth, and carry its state beside the support
of every other hypothesis. The search then becomes a nondeterministic automaton over windows. Its states have two
forms. $U(S)$, for a non-empty $S \subseteq Q^+$, means that no origin has been chosen yet and that $S$ is the set of
states the cloud occupies. $V(S, q)$, with $q \in Q^+$ and $S \subseteq Q^+ \setminus \{q\}$, means that the chosen
origin sits at $q$ and that the pairs carrying any other origin occupy exactly the states of $S$. The initial state
is $U(Q^+)$. The accepting states are the $V(\varnothing, q)$.

Take a byte $b$ at $U(S)$. The rename fires when $q_0 \in S$ and $q_0$ is not re-entrant. Let $T$ be the live image
of $S$ under $b$, with the contribution of $q_0$ set aside when the rename fires. A fresh origin is born at $r =
\delta(q_0, b)$ when $r$ is live and either $S$ contains an accepting state or the rename fired. Without a birth the
successor is $U(T)$, when $T$ is non-empty. With a birth there are two successors. Declining the birth moves to $U(T
\cup \{r\})$. Choosing it moves to $V(T, r)$, which is allowed only when $r \notin T$. Now take a byte $b$ at $V(S,
q)$. Let $q' = \delta(q, b)$; there is no successor when $q'$ is not live. Let $T$ be the live image of $S$,
together with $\delta(q_0, b)$ when that state is live and $S \cup \{q\}$ contains an accepting state. There is no
successor when $q' \in T$. Otherwise the successor is $V(T, q')$. No rename arises in a $V$ state: a non-re-entrant
$q_0$ is unoccupied after the first byte, and a re-entrant one disables the rename.

\begin{theorem}[Chosen-origin automaton]
\label{thm:chosen}
The automaton accepts exactly the windows the model certifies. It is unambiguous: a certified window has exactly one
accepting run, which chooses the window's origin at that origin's birth. It has at most
\[ K \;=\; 2^n + n\, 2^{n-1} \;=\; (n + 2)\, 2^{n-1} \]
states, and at most two successors per state and byte.
\end{theorem}

\begin{proof}
Two invariants hold along every run, by induction on the bytes read. In $U(S)$ after a word $u$, the set $S$ is the
support of $C_u$. The gate reads only which occupied states accept. The direct step maps the support to its live
image. The rename is read from $q_0 \in S$ as in Theorem~\ref{thm:three}, and a fresh origin occupies $\delta(q_0,
b)$. In $V(S, q)$ after $u$, the cloud $C_u$ contains a pair $(q, o)$, where $o$ is the origin chosen when the run
left $U$, and $S$ is the set of states occupied by pairs with an origin other than $o$. The pair $(q, o)$ moves to
$(\delta(q, b), o)$ when that state is live. Otherwise $o$ dies, and no later cloud contains it, since every origin
born later carries a larger offset. The other pairs move to the live image of $S$. The fresh origin, when born, is
another origin and lands at $\delta(q_0, b)$, and the gate reads the support $S \cup \{q\}$ of $C_u$. So the
successor $V(T, q')$ records exactly the state of $o$ and the support of the others, as long as $q' \notin T$. When
$q' \in T$ the run stops, and nothing is lost. The state $q'$ then carries $o$ and a second origin. By determinism
the two move together or die together at every later step, so no later cloud has $o$ as its only origin.

A run in $V(\varnothing, q)$ after $w$ thus says $C_w = \{(q, o)\}$ with $o$ in the window, which is certification.
Conversely let $w$ be certified with origin $o$. At the byte $w_o$ the origin $o$ is born at $\delta(q_0, w_o)$, by
the seed or by the rename. That state carries no other origin, for otherwise $o$ could never become the only one.
Choose $o$ there and decline every other birth. The chosen pair survives to the end, since a pair carrying $o$ is in
$C_w$. It never shares a state with another origin, since both would then reach $C_w$. Its final $S$ is empty,
because every pair of $C_w$ carries $o$. For unambiguity, an accepting run chooses the origin that is alone at the
end. That origin is the unique origin of $C_w$, and it has one birth. Before the choice the run is deterministic,
and after it the transitions are deterministic. For the count, there are $2^n$ sets $S$ for the $U$ states, counting
the empty support, which is never entered. For the $V$ states there are $n$ choices of $q$, and $S$ is a subset of
the remaining $n - 1$ states. Only a birth before the choice branches.
\end{proof}

\begin{corollary}[Shortest window]
\label{cor:chosen}
Breadth-first search over the automaton decides whether the model certifies any window and returns a shortest one.
It takes $O(256 \cdot n \cdot K)$ time and $O(n \cdot K)$ space. The origin is the depth at which the returned path
leaves $U$. If a certified window exists, one of length at most $K - 1$ exists. Existence of a model certificate is in
PSPACE in the size of the compiled automaton.
\end{corollary}

\begin{proof}
A path from $U(Q^+)$ to an accepting state spells a certified window, and conversely. Breadth-first order finds a
shortest one, and the origin is read off the one $U$-to-$V$ edge of the path. A shortest accepting path repeats no
state, so its length is below $K$. A state is written in $O(n)$ bits, and a successor is computed in polynomial
time. A path of length below $K$ is guessed holding one state and a counter of $O(n)$ bits. That places existence in
nondeterministic polynomial space, which is PSPACE.
\end{proof}

\begin{remark}
\label{rem:determinize}
Determinizing the chosen-origin automaton gives back the three-status quotient. After a word with three-status key
$(S, B)$, the set of automaton states reached is $U(S \cup B)$ together with $V((S \cup B) \setminus \{q\},\, q)$ for
every $q \in S$. A $V$ run survives exactly while its origin is alone at its state, which is the single status, and
its other support is then the rest of the occupied states. The subset construction stays within these special
subsets, which is why no second exponential appears.
\end{remark}

Release 2.2 of munch is the release that ships the search of Corollary~\ref{cor:chosen}. It is named
\code{shortest\_split\_window()} there, beside the window decision \code{is\_split\_window()} of release 1.4.0. The
search holds one node per automaton state reached, packed as a chosen state and a bit per live state. It reads the
window and its origin back along the trail of parents. A node budget, $2^{20}$ by default, bounds the nodes held.
Running out is reported as such, never as an answer. The library's test reproduces the first six minima of the family of
Section~\ref{sec:long}. The figure program of Section~\ref{sec:evaluation} runs the same search over the named token
sets, and Table~\ref{tab:windows} reports the nodes it holds when its shortest window appears.

\subsection{PSPACE-completeness}
\label{sec:complexity}

The upper bound of Corollary~\ref{cor:chosen} is matched by a lower bound. The lower bound holds with one token kind
and five symbols. The source problem is intersection non-emptiness: given complete DFAs $A_1, \ldots, A_m$ over a
common alphabet, is some word accepted by all of them? Kozen proved it PSPACE-complete~\cite{kozen1977lower}. We take
the source DFAs over the alphabet $\{\code{a}, \code{b}\}$, which loses nothing. A fixed-length binary code of the
letters, decoded by a polynomial number of added states with a rejecting sink for malformed codes, carries any
alphabet to this one. We also keep in each $A_i$ only the states reachable from its start, which changes no language.

\begin{theorem}[PSPACE-completeness]
\label{thm:pspace}
Deciding whether a compiled token set admits a model certificate of any length is PSPACE-complete. So is deciding
whether it admits a witnessed one. Both remain PSPACE-hard for token sets with one token kind over five symbols,
whose start state neither accepts nor is re-entered and whose every other state accepts.
\end{theorem}

Membership for the plain problem is Corollary~\ref{cor:chosen}. For the witnessed problem it is
Theorem~\ref{thm:witnessed} below. The rest of this section is the reduction. Let $A_i = (Q_i, \{\code{a},
\code{b}\}, \delta_i, s_i, F_i)$ for $i = 0, \ldots, m - 1$, with $m \geq 1$. The token DFA has the symbols
$\code{a}, \code{b}, \code{r}, \code{t}, \code{\#}$. Its states are the root $q_0$, an error state $E$, a guard $G$,
a result state $Y$, timers $T_0, \ldots, T_m$, headers $H_0, \ldots, H_m$, and a disjoint copy $(i, q)$ of every
state $q$ of every $A_i$. The root is the only state that does not accept, and nothing enters it. Every other state
accepts the one token kind. The root has the two transitions $q_0 \xrightarrow{\code{r}} H_0$ and $q_0
\xrightarrow{\code{\#}} Y$. The other transitions are in Table~\ref{tab:reduction}. An entry marked undefined is
absent, and on the remaining bytes of the alphabet no state has a transition.

\begin{table}
\centering
\caption{The transitions of the reduction's token DFA at every state but the root, which has
$q_0 \xrightarrow{\code{r}} H_0$ and $q_0 \xrightarrow{\code{\#}} Y$ only. The letter $x$ ranges over the source
alphabet $\{\code{a}, \code{b}\}$.}
\label{tab:reduction}
\begin{tabular}{lcccc}
\toprule
State & \code{r} & \code{t} & \code{\#} & $x$ \\
\midrule
$E$ & $T_m$ & $E$ & $E$ & $E$ \\
$G$ & $T_m$ & $E$ & undefined & $G$ \\
$Y$ & $T_m$ & $E$ & $E$ & $E$ \\
$T_j$, $j > 0$ & $T_{j-1}$ & $E$ & $E$ & $E$ \\
$T_0$ & $T_0$ & $G$ & $E$ & $E$ \\
$H_j$, $j < m$ & $H_{j+1}$ & $(j, s_j)$ & $E$ & $E$ \\
$H_m$ & $H_m$ & $G$ & $E$ & $E$ \\
$(i, q)$ & $T_m$ & $E$ & undefined if $q \in F_i$, else $E$ & $(i, \delta_i(q, x))$ \\
\bottomrule
\end{tabular}
\end{table}

The construction has $\sum_i |Q_i| + 2m + 6$ states and takes polynomial time. Every state is live. The root reaches
$H_j$ on $\code{r}^{j+1}$, $Y$ on \code{\#}, $E$ on \code{\#\#}, $T_m$ on \code{\#r} and each lower timer on further
\code{r}'s. It reaches $G$ on $\code{r}^{m+1}\code{t}$, the start $(i, s_i)$ on $\code{r}^{i+1}\code{t}$, and every
other $(i, q)$ on the word that reaches $q$ in $A_i$ from there. Every state but the root accepts, and the root has
an accepting successor, so all are co-accessible. The root is not re-entrant, since no transition targets it, and
the rename is in force. Two facts about the table carry the proof. First, on $\code{r}, \code{t}, \code{a},
\code{b}$ every state but the root has a transition. So on these letters no pair at a non-root state ever dies.
Second, every occupied state after the first byte is a non-root state and accepts. So the gate is open at every step
after the first, and it is open at the first because $C_0$ holds accepting states.

\begin{lemma}[Common word to window]
\label{lem:reduction-forward}
If every $A_i$ accepts $w$, then $W = \code{r}^{m+1}\, \code{t}\, w\, \code{\#}$ is certified at origin $|W| - 1$.
Moreover $W$ is itself completely tokenizable, with its final byte a token of its own.
\end{lemma}

\begin{proof}
Follow the cloud. During the $m + 1$ bytes \code{r}, every pair at $E$, $G$, $Y$, a timer or a source copy moves
into the timers and down to $T_0$, where it stays. Every pair at a header moves up to $H_m$ and stays. The root's
pair is renamed to origin $0$ at $H_0$. The gate is open at every byte, so a fresh origin $j$ is born at $H_0$ at
each byte $j = 0, \ldots, m$; the rename and the seed coincide at $j = 0$. After the block, origin $j$ sits at
$H_{m-j}$, one origin per header, with $(T_0, \before)$ and $(H_m, \before)$ beside them. The byte \code{t} sends
$T_0$ and $H_m$ to $G$, and $H_i$ to $(i, s_i)$ for $i < m$. The root has no \code{t} edge, so no origin is born.
The letters of $w$ keep $G$ fixed and drive each $(i, s_i)$ to $(i, \delta_i(s_i, w))$, with no birth. On the final
\code{\#} every pair at $G$ dies. Every source pair dies too, because its state is in $F_i$. The open gate seeds
$(Y, |W| - 1)$, and the final cloud is that singleton. For the scan, the run of $W$ from the root passes $H_0,
\ldots, H_m$, enters $G$ on \code{t} and stays there through $w$, accepting at every byte. $G$ has no \code{\#}
edge, so the first token is $W$ without its last byte. The scan restarts at \code{\#}, which the root consumes into
$Y$.
\end{proof}

\begin{lemma}[Window to common word]
\label{lem:reduction-back}
Suppose the model certifies some window. Then it certifies a prefix of it of the form $P\, \code{r}^{L}\, \code{t}\,
w\, \code{\#}$, with $L \geq m + 1$ and $w \in \{\code{a}, \code{b}\}^*$ accepted by every $A_i$.
\end{lemma}

\begin{proof}
Take a certified window and cut it at the first byte after which the cloud is certified. Call the result $W' = w_0
\cdots w_{k-1}$, so $C_k$ is certified and $C_{k-1}$ is not. A byte outside the five symbols empties the cloud, which is
absorbing and never certified. So every byte of $W'$ is one of the five, and $C_{k-1}$ is non-empty. The last byte is
\code{\#}. To see this, suppose it were one of $\code{r}, \code{t}, \code{a}, \code{b}$. Every pair of $C_{k-1}$ at a
non-root state then survives with its origin intact, since the rename touches only the root. For $k = 1$ the surviving
pairs carry $\before$. For $k > 1$ they carry whatever made $C_{k-1}$ uncertified: $\before$, or two distinct origins. A
fresh origin only adds a pair, and cannot certify an uncertified cloud. So $C_k$ would not be certified. On \code{\#},
every pair moves to $E$ or dies, and the open gate seeds $(Y, k - 1)$. Certification at origin $k - 1$ therefore needs
every pair of $C_{k-1}$ to die. So $C_{k-1}$ occupies only $G$ and accepting source copies.

Now choose a sentinel. If some byte of $w_0 \cdots w_{k-2}$ is \code{\#}, let $j$ be the position of the last one.
The cloud $C_j$ is non-empty, since the empty cloud is absorbing and $C_{k-1}$ is not empty. The gate is open, so
$(Y, j) \in C_{j+1}$. Otherwise take $(E, \before) \in C_0$ and set $j = -1$. The bytes $w_{j+1}, \ldots, w_{k-2}$
are not \code{\#}, so they are among $\code{r}, \code{t}, \code{a}, \code{b}$, on which no pair at a non-root
state dies. So the sentinel
survives to $C_{k-1}$ with its origin, following the table from $Y$ or $E$. From $Y$, $E$, the timers and $G$, these
four letters lead only to $Y$, $E$, the timers and $G$. So the sentinel never reaches a header or a source copy. To
die on the final \code{\#} it must sit at $G$ in $C_{k-1}$. Call $z = w_{j+1} \cdots w_{k-2}$ the word the sentinel
reads. Among these states, $G$ is entered only from $T_0$ on \code{t}; it is left on \code{r} and \code{t}, and
kept on \code{a} and \code{b}.
So the last \code{t} of $z$ is read from $T_0$, and only the letters \code{a} and \code{b} follow it. That is, $z =
z' \code{t} w$ with $w \in \{\code{a}, \code{b}\}^*$, and the sentinel is at $T_0$ after $z'$. $T_0$ is entered
only from $T_1$ or from itself, on \code{r}. Every other letter from a timer leads to $E$ or $G$. Let $\code{r}^L$
be the longest run of \code{r}'s ending $z'$. The state before that run is reached by a letter other than \code{r},
or is the sentinel's initial state. So it is among $Y$, $E$, $G$. If $L = 0$ the sentinel is there after $z'$, not
at $T_0$. Otherwise the run enters $T_m$ on its first byte and needs $m$ more to reach $T_0$: $L \geq m + 1$. So $W'
= P\, \code{r}^{L}\, \code{t}\, w\, \code{\#}$.

During the last $m + 1$ bytes of the run $\code{r}^L$ the gate is open, and the root's \code{r} edge is live. So an
origin is born at $H_0$ at each of them. The one born $i$ bytes before the run's end sits at $H_i$ when the run ends,
for $i = 0, \ldots, m$. These origins are distinct, one per header. The \code{t} sends $H_i$ to $(i, s_i)$ for $i <
m$, and $w$ drives each to $(i, \delta_i(s_i, w))$. Every pair of $C_{k-1}$ dies on the final \code{\#}, and a
source copy dies there exactly when its state is in $F_i$. So $\delta_i(s_i, w) \in F_i$ for every $i$.
\end{proof}

\begin{proof}[Proof of Theorem~\ref{thm:pspace}]
Lemma~\ref{lem:reduction-forward} maps a common word to a certified window. That window is witnessed, since $W$ is
its own completely tokenizable occurrence. Lemma~\ref{lem:reduction-back} maps any certified window, witnessed or
not, to a common word. The construction is polynomial, so both problems are PSPACE-hard. They are in PSPACE by
Corollary~\ref{cor:chosen} and Theorem~\ref{thm:witnessed}.
\end{proof}

The reduction fixes the minimum length exactly, at the source's plus $m + 3$.

\begin{corollary}[Minimum length]
\label{cor:reduction-length}
Let $\ell$ be the length of a shortest word every $A_i$ accepts. Then a shortest model-certified window of the
constructed token set has length $\ell + m + 3$, and so does a shortest witnessed one.
\end{corollary}

\begin{proof}
Lemma~\ref{lem:reduction-forward} gives a witnessed window of length $\ell + m + 3$. By
Lemma~\ref{lem:reduction-back}, any certified window has a certified prefix ending in $\code{r}^{L}\, \code{t}\, w\,
\code{\#}$ with $L \geq m + 1$ and $|w| \geq \ell$. So its length is at least $\ell + m + 3$. A witnessed window is
certified, so the same bound holds for it.
\end{proof}

The theorem concerns the model. It says nothing about the complexity of deciding Definition~\ref{def:window} itself,
which the model only approximates from below. It also concerns the compiled automaton. A token set given as regular
expressions may compile to exponentially many states, and no bound in the expression size is claimed. What it rules
out is a polynomial algorithm for the model's existence question, unless $\mathrm{P} = \mathrm{PSPACE}$. The
exponential worst case of Theorem~\ref{thm:decision} is therefore not an artifact of the quotient.

\subsection{Shortest windows can be long}
\label{sec:long}

Corollary~\ref{cor:chosen} bounds a shortest certified window by $K - 1$, which is exponential in $n$. The bound is
not far off. The reduction instantiated on cycles gives windows of length $2^{\Omega(\sqrt n)}$, and the figure
program searches their exact minima.

For $m \geq 1$ let the source $A_i$, for $i = 1, \ldots, m$, be a cycle of $i$ states numbered $0, \ldots, i - 1$,
placed behind the header $H_{i-1}$. It starts at $0$, both letters advance by one modulo $i$, and it accepts at
state $i - 1$ alone. It accepts exactly
the words whose length is $-1$ modulo $i$. Let $L_m = \mathrm{lcm}(1, \ldots, m)$. A word every source accepts has
length $-1$ modulo $L_m$, and $\code{a}^{L_m - 1}$ is a shortest one. The reduction gives a token set with
\[ n_m \;=\; \sum_{i=1}^{m} i + 2m + 6 \;=\; \frac{m^2 + 5m + 12}{2} \]
states. By Corollary~\ref{cor:reduction-length} its shortest certified window, plain or witnessed, has length $R_m =
L_m + m + 2$. One such window is $\code{r}^{m+1}\, \code{t}\, \code{a}^{L_m - 1}\, \code{\#}$. It is certified at
its last byte and completely tokenizable as two tokens.

\begin{theorem}[Long windows]
\label{thm:long}
For every $m \geq 1$ the token set above has $n_m$ live states and a shortest certified window of length $R_m \geq
2^m / (m + 1) + m + 2$. Hence $R_m = 2^{\Omega(\sqrt{n_m})}$. No polynomial $p$ bounds the length of a shortest
model-certified window by $p(n)$ over all compiled token sets, nor of a shortest witnessed one.
\end{theorem}

\begin{proof}
It remains to show $L_m \geq 2^m / (m + 1)$. Fix a prime $p$ and write $v_p$ for the exponent of $p$ in an integer.
By the formula for the exponent of $p$ in a factorial,
\[ v_p\binom{m}{k} \;=\; \sum_{j \geq 1} \left( \left\lfloor \frac{m}{p^j} \right\rfloor - \left\lfloor
\frac{k}{p^j} \right\rfloor - \left\lfloor \frac{m - k}{p^j} \right\rfloor \right). \]
Each term is $0$ or $1$, since $m = k + (m - k)$ and the floor of a sum exceeds the sum of the floors by at most one.
Terms with $p^j > m$ vanish, so $v_p\binom{m}{k} \leq \lfloor \log_p m \rfloor = v_p(L_m)$. This holds for every
prime, so every $\binom{m}{k}$ divides $L_m$. In particular $L_m \geq \max_k \binom{m}{k} \geq 2^m / (m + 1)$, since
the $m + 1$ binomial coefficients sum to $2^m$. Since $n_m$ is quadratic in $m$, $m = \Theta(\sqrt{n_m})$ and the
growth follows. For the last claim, $R_m / n_m^d$ is unbounded for every fixed degree $d$.
\end{proof}

The family's sizes are not every $n$, but they are dense enough. Since $n_{m+1} - n_m = m + 3 \leq n_m$, for every
$n \geq 9$ the member with the largest $n_m \leq n$ has $n < 2 n_m$. So some token set with at most $n$ live states
has a shortest window of length $2^{\Omega(\sqrt n)}$. Table~\ref{tab:family} gives the exact minima for $m = 1,
\ldots, 9$, searched by the chosen-origin search of the figure program. The program also replays each window
through the plain cloud at the claimed origin and scans it. The result is explicit: a shortest window is long, not
merely hard to find. It does not exclude a compressed description of long windows. And it concerns the model; the
semantic question meets the constraint only through the model.

% Every figure of the table is printed and asserted by paper/figures/window_search.cpp, the family section.
\begin{table}
\centering
\caption{The long-window family: live states, $\mathrm{lcm}(1, \ldots, m)$, the shortest certified window's length,
equal to $L_m + m + 2$ on every row, and the nodes the chosen-origin search holds when it appears. Every window is
certified at its last byte and scans as two tokens.}
\label{tab:family}
\begin{tabular}{crrrr}
\toprule
$m$ & States $n_m$ & $L_m$ & Shortest window & Nodes held \\
\midrule
1 & 9 & 1 & 4 & 18 \\
2 & 13 & 2 & 6 & 41 \\
3 & 18 & 6 & 11 & 110 \\
4 & 24 & 12 & 18 & 253 \\
5 & 31 & 60 & 67 & 1{,}226 \\
6 & 39 & 60 & 68 & 1{,}781 \\
7 & 48 & 420 & 429 & 11{,}038 \\
8 & 58 & 840 & 850 & 27{,}017 \\
9 & 69 & 2{,}520 & 2{,}531 & 87{,}518 \\
\bottomrule
\end{tabular}
\end{table}

\section{Specializations}
\label{sec:specialization}

\subsection{Length one}
\label{sec:length-one}

At length one the model collapses to the published certificate. The correspondence is exact at the level of the
\emph{shipped predicate}, the one that withholds vacuously certified bytes, rather than of the bare condition.
Throughout this section the token set is non-empty, read through its positive-width equivalent where a token matches the
empty string, and $q_0$ is live.

Call a byte $b$ \emph{useful} when $\delta(q_0, b)$ is defined and live. The predicate \code{is\_split\_point($b$)}
of~\cite{nidhogg2026splitpoints} holds exactly when $b$ is useful, no live state other than $q_0$ has a $b$-transition
into a live state, and $q_0$ is not re-entrant in $A^+$. Two quantifier readings coincide here without loss. A
transition target that is co-accessible makes its source co-accessible, so ``reachable state with a live-target
$b$-transition'' and ``live state with a live-target $b$-transition'' name the same states.

\begin{theorem}[Specialization]
\label{thm:specialization}
For every byte $b$, the model certifies $(b, 0)$, the only origin a length-one window admits, if and only if
\code{is\_split\_point($b$)} holds.
\end{theorem}

\begin{proof}
Compute $C_1$ explicitly. $C_0$ is every live state paired with $\before$. The initial cloud contains a live accepting
state, since $q_0$ is live and an accepting state reachable from it is trivially co-accessible, so the seed's acceptance
gate is open at offset 0. Reading $b$ therefore yields exactly:

\begin{itemize}
\item
from the seed, the pair $(\delta(q_0, b),\, 0)$, present if and only if $b$ is useful;
\item
from the direct step on $(q_0, \before)$, present if and only if $b$ is useful: the pair
$(\delta(q_0, b),\, 0)$ when $q_0$ is not re-entrant, by the rename, and $(\delta(q_0, b),\, \before)$ when it is;
\item
from the direct step on every other live state $q$ that has a live-target $b$-transition, the
pair $(\delta(q, b),\, \before)$.
\end{itemize}

If \code{is\_split\_point($b$)} holds, only the first two items contribute and the rename fires. In that case both
surviving contributions are the single pair $(\delta(q_0, b),\, 0)$, so $C_1$ is the singleton and the model certifies
$(b, 0)$. Conversely, suppose the model certifies $(b, 0)$, so $C_1$ is non-empty and unanimous at origin $0$. If $b$
were not useful, the first two items would be absent and every pair of $C_1$ would carry $\before$ from the third, or
$C_1$ would be empty. Either way certification fails, so $b$ is useful. If some live $q \neq q_0$ had a live-target
$b$-transition, the third item would place a $\before$-pair in $C_1$, breaking unanimity; so none does. If $q_0$ were
re-entrant, the second item would place $(\delta(q_0, b),\, \before)$ in $C_1$ beside the seed's origin-$0$ pair,
breaking unanimity; so $q_0$ is not re-entrant. These are exactly the predicate's three clauses.
\end{proof}

The vacuous case lands on the predicate's side of the published condition-versus-predicate distinction by construction.
A byte no live state consumes empties the cloud, which the model refuses, exactly as the shipped predicate withholds a
byte the bare condition certifies vacuously. At length one the model is therefore exactly the shipped predicate. The
multi-byte construction is the predicate's conservative continuation. The artifact additionally asserts the agreement
byte for byte on every grammar of the evaluation before each search. After Theorem~\ref{thm:specialization} that check
verifies the implementation rather than the claim.

\subsection{The prefix-code case}
\label{sec:codes}

The other specialization is by token set rather than by length. A \emph{prefix code} $X \subseteq \Sigma^+$ has no word
that is a proper prefix of another. Read as a token set, a prefix code is the setting where the certificate needs
neither maximal-munch priority nor the cloud's origin bookkeeping. At any boundary at most one codeword begins, so the
completely tokenizable inputs are exactly $X^*$, each with the one factorization the scan computes. Code theory has two
notions of a boundary forced by bounded context. A \emph{synchronizing pair} of $X$~\cite{berstel2010codes} is a pair
$(x, y) \in X^* \times X^*$ such that $u x y v \in X^*$ implies $u x \in X^*$ and $y v \in X^*$ for all $u, v \in
\Sigma^*$. In the factor formulation of Fici, Romana, Sciortino and Urbina~\cite{fici2025morphisms}, a split $w = w_1
w_2$ of a factor $w$ of $X^*$ is \emph{synchronizing} when every occurrence of $w$ in a word of $X^*$ has a
factorization boundary between $w_1$ and $w_2$. Neither is a certified window as it stands, since neither says which
codeword covers the window's final byte, and the certificate says nothing about the halves being in $X^*$. The exact
relation is the following.

\begin{proposition}[Prefix codes]
\label{prop:codes}
Let $X$ be a prefix code read as a token set, and let $W$ be a window occurring in some word of $X^*$, with
origin $o$, split as $W = W_{<o}\, W_{\geq o}$ at the origin.
\begin{enumerate}
\item $(W, o)$ is certified if and only if the split is synchronizing and $W_{\geq o}$ is a prefix of some codeword.
\item If $(x, y)$ is a synchronizing pair with $y$ nonempty and $y = c_1 \cdots c_m$ is its factorization into
codewords, then
$(xy,\ |x| + |c_1 \cdots c_{m-1}|)$ is certified. Conversely, if $(W, o)$ is certified with $W_{<o} \in X^*$ and
$W_{\geq o} \in X$, then $(W_{<o}, W_{\geq o})$ is a synchronizing pair.
\end{enumerate}
\end{proposition}

\begin{proof}
Throughout, a completely tokenizable input is a word of $X^*$ and its token boundaries are the boundaries of its one
factorization, since at a boundary the codeword the text spells is the only codeword that begins there.

(1) If $(W, o)$ is certified, every occurrence has a boundary at offset $o$, so the split is synchronizing. By
certification, the codeword covering the final byte begins at $o$, so $W_{\geq o}$ lies within that codeword from its
start and is a prefix of it. Conversely, take an occurrence of $W$ at offset $t$ in a word of $X^*$. The split being
synchronizing puts a boundary at $t + o$, where some codeword $c'$ begins. Let $c$ be a codeword with prefix $W_{\geq
o}$. Were $c'$ to end inside $W_{\geq o}$, it would be a proper prefix of $c$, which a prefix code forbids; so $c'$
extends through $W_{\geq o}$, and the token covering the final byte begins at $t + o$.

(2) Let $(x, y)$ be a synchronizing pair and take an occurrence $u x y v \in X^*$. The membership $u x \in X^*$ places a
boundary before $y$. From that boundary, $y v \in X^*$ with $y \in X^*$ factors through $y$'s own codewords, so $c_m$
begins at $|u x| + |c_1 \cdots c_{m-1}|$ and covers the final byte of $xy$. The factorization follows $y$'s codewords
because at each of their starts the codeword the text spells is the one $y$ contains. Conversely, if $(W, o)$ is
certified with $W_{<o} \in X^*$ and $W_{\geq o} \in X$, then every $u W v \in X^*$ has a boundary at $|u| + o$, so $u
W_{<o}$ is a concatenation of codewords and so is $W_{\geq o} v$, which is the synchronizing-pair condition.
\end{proof}

The proposition places the classical case in three parts. Over a prefix code the certified windows that occur are the
synchronizing splits whose right half sits inside one codeword. A synchronizing pair of Berstel, Perrin and Reutenauer
with a nonempty right component is a certified window on its concatenation, the origin at the start of the right
component's last codeword. A certified window whose left half is in $X^*$ and whose right half is a codeword is a
synchronizing pair. The occurrence hypothesis sets the vacuous case of Definition~\ref{def:window} aside: a window no
word of $X^*$ contains is certified whatever its right half spells. The named token sets of Section~\ref{sec:evaluation}
fall outside the code case. An identifier is a proper prefix of a longer identifier, so no C-like token set is a prefix
code. The JSON number \code{1} is a proper prefix of \code{11}, so JSON is not one either. Over $\{\code{a}, \code{ab},
\code{b}\}$ the word \code{ab} has two factorizations, and maximal munch, not the code, picks one. The random sample of
Section~\ref{sec:evaluation} is drawn without that restriction and contains prefix codes, which the proposition covers.
Beyond the code case the certificate is defined through the scan, the origin bookkeeping of Section~\ref{sec:model} is
what decides it, and that reach is what the evaluation measures.

\section{Strictness of the model}
\label{sec:strictness}

The model refuses windows a greedy scanner would allow. The conservatism is deliberate: the seed rule uses acceptance as
the only license a token needs to begin, which over-approximates greedy behaviour by design. This section exhibits one
source of conservatism and shows that nonvacuous strictness begins at length two. Both witnesses below are asserted
artifact rows. The assertion checks that the model refuses the window. The same assertion checks that an exhaustive
oracle over every completely tokenizable input up to a length bound finds, at every occurrence, the token covering the
window's final byte beginning at the claimed origin. The occurrence count is pinned exactly. The covering-token check is
the property of Definition~\ref{def:window}. Checking merely that some token begins at the origin passes false
covering-origin witnesses, such as $\{\code{a}, \code{abx}, \code{b}, \code{x}\}$ at \code{ab}, where the input
\code{ab} tokenizes as \code{a}$\mid$\code{b} while the fixed boundary at the occurrence remains a safe cut.

\emph{Witness one.} Over $\{\code{a}, \code{ab}, \code{b}\}$ the window \code{ab} is semantically certified at origin 0,
and universally so, not merely to the oracle's bound. The byte \code{a} occurs only as a token's first byte, so a token
begins at every occurrence offset $t$. Maximal munch there prefers \code{ab} over \code{a}, so the covering token of the
final byte begins at $t$. The oracle confirms the argument over all inputs to length 14, with 98{,}305 occurrences and
zero violations. The model refuses the window. After \code{a} the cloud is the single pair carrying origin 0, but
\code{a} is a token, so reading \code{b} seeds a competing trajectory at origin 1. The competing trajectory is the
segmentation \code{a}$\mid$\code{b} that greedy scanning never chooses, and unanimity is lost.
Section~\ref{sec:factorization} shows that such segmentations are the whole of the surplus.

\emph{Witness two.} Over $\{\code{ab}, \code{abc}, \code{c}\}$ the window \code{abc} is semantically certified at origin
0, again universally. The byte \code{a} occurs only token-initially, so a token begins at $t$, and maximal munch prefers
\code{abc} over \code{ab} there. The oracle confirms it over all inputs to length 12 with 932 occurrences and zero
violations. Both witnesses instantiate the same $(u,\, uv,\, v)$ shape. The two differ in prefix depth rather than
mechanism, the competing origin arriving one byte in for the first and two bytes in for the second. In the second the
accepting proper prefix \code{ab} seeds \code{c}, the segmentation \code{ab}$\mid$\code{c} that maximal munch forgoes.

Both witnesses have length at least two, and that is not an accident of the examples.

\begin{corollary}[Non-vacuous strictness begins at length two]
\label{cor:strictness}
Let $b$ be a byte occurring in some completely tokenizable input. If the model refuses $(b, 0)$, then $(b, 0)$ is
not a certified split window. The occurrence hypothesis is necessary: a byte no completely tokenizable input
contains satisfies Definition~\ref{def:window} vacuously while the model refuses its emptied cloud, and such
vacuous disagreements are not strictness.
\end{corollary}

\begin{proof}
Since $b$ occurs, the final segmentation's covering token consumes $b$ at that occurrence, a transition from a live
state into a live state. Given that transition, if only $q_0$ has a live-target $b$-transition and $q_0$ is not
re-entrant, the predicate reports $b$. In that case the model certifies $(b, 0)$ by Theorem~\ref{thm:specialization},
contrary to assumption, so the exact condition of~\cite{nidhogg2026splitpoints} fails for $b$. The condition's necessity
theorem constructs a completely tokenizable input placing an occurrence of $b$ strictly inside a token. At that
occurrence the token containing $b$ begins before $b$, so $(b, 0)$ is not certified.
\end{proof}

Non-vacuous conservatism is therefore a strictly multi-byte phenomenon: at length one the model is exact for occurring
bytes, and the shortest strict refusals have length two, a bound witness one attains. The $\{\code{a}, \code{abx},
\code{b}, \code{x}\}$ family of Section~\ref{sec:preliminaries} plays the opposite role in the artifact, a negative
control for the covering-origin property itself. Cutting at the family's fixed boundary is safe, since \code{a} occurs
only token-initially. The covering origin genuinely varies, and the artifact pins those violations exactly. An oracle
that misses them has lost its teeth. Negatives in every table of this paper are claims about the model, never about the
language.

\section{What the cloud represents}
\label{sec:factorization}

The cloud was built as a conservative over-approximation, and Section~\ref{sec:strictness} exhibited the surplus. This
section says exactly what the surplus is. Let $X$ be the token language, the set of non-empty words $A$ accepts from
$q_0$. Let $\mathcal{G}$ be the set of completely tokenizable inputs, so $\mathcal{G} \subseteq X^*$. A
\emph{factorization} of a word $x \in X^*$ is a sequence of words of $X$ whose concatenation is $x$, with no greedy
choice imposed. The cloud is exactly the set of histories of all factorizations of all words of $X^*$ around the
window. Everything else in this section follows from that. It explains why the vacuous window of
Section~\ref{sec:preliminaries} is certified. It locates a class where the model is exact. It finds a shortest
witnessed window exactly. And it decides whether the model misses any witnessed certificate at all.

\subsection{Factorization semantics}
\label{sec:semantics}

Fix a window $W = w_0 \cdots w_{k-1}$. Take a word $x = u W v \in X^*$ and a factorization of $x$. Let $s$ be the
start of the factor covering the byte at $|u| + k - 1$. Let $q = \delta^*(q_0,\, x[s \,..\, |u| + k))$ be that
factor's state after the window's last byte. Let $\omega$ be $\before$ when $s < |u|$, and $s - |u|$ otherwise. The
set $F(W)$ collects all pairs $(q, \omega)$ so obtained, over every $u$, $v$ and factorization.

\begin{theorem}[Factorization semantics]
\label{thm:factorization}
For every non-empty window, $C_k = F(W)$. Every pair of $C_k$ is realized by some $x = u W v$ with $|u| \leq n$ and
$|v| \leq n - 1$. In particular $C_k$ is non-empty exactly when $W$ is a factor of a word of $X^*$.
\end{theorem}

\begin{proof}
First $F(W) \subseteq C_k$. Fix $x = u W v$ and a factorization. For $j \in 1..k$ let $s_j$ be the start of the
factor covering the byte at $|u| + j - 1$, with $\rho_j$ and $\omega_j$ defined from it as in
Section~\ref{sec:soundness}. The argument of Lemma~\ref{lem:representation} applies with ``factor'' in place of
``token of the final segmentation''. It used only that tokens are accepted words placed end to end. Each $\rho_j$ is
live, reached by a prefix of an accepted factor and completed by its rest. If $s_1 < |u|$, the state $p =
\delta^*(q_0, x[s_1 \,..\, |u|))$ is live and $(p, \before) \in C_0$. The direct step gives $(\rho_1, \before)$. The
rename cannot fire, since $p$ has consumed a byte and $p = q_0$ only when $q_0$ is re-entrant. If $s_1 = |u|$, the
seed fires, because $C_0$ holds an accepting state, $X$ being non-empty, and it gives $(\rho_1, 0)$. At a continuing
byte the direct step carries the pair with its origin, and again the rename cannot fire. At a byte beginning a factor
the previous factor $x[s_j \,..\, |u| + j)$ is in $X$. So $\rho_j$ accepts and the seed gives $(\rho_{j+1}, j)$.

Now $C_k \subseteq F(W)$, by induction on $j$. The claim is that every pair of $C_j$ is the pair of some $x = u\,
W[0 \,..\, j)\, v$ with a factorization, where $|u| \leq n$ and $v$ is a shortest accepting completion of the
covering factor's state, so that $|v| \leq n - 1$. For $j = 1$ the pairs of $C_1$ arise three ways. A direct step
from $(p, \before)$ with $p \neq q_0$ gives $(\delta(p, w_0), \before)$. Let $u$ be a shortest word with
$\delta^*(q_0, u) = p$; it is non-empty and of length at most $n - 1$. Let $v$ be a shortest accepting completion
from $\delta(p, w_0)$, which exists since that state is live. Then $u w_0 v \in X$ is one factor, starting before the
window. A direct step from $(q_0, \before)$ that keeps $\before$ happens only when $q_0$ is re-entrant. Then a
shortest non-empty $u$ with $\delta^*(q_0, u) = q_0$ has length at most $n$, and the same completion gives one factor
$u w_0 v$. The fresh origin, from the seed or from the rename of $(q_0, \before)$, gives $(\delta(q_0, w_0), 0)$. It
is realized with $u$ empty by the factor $w_0 v$ beginning at the window.

For the step, a pair of $C_{j+1}$ is either a direct successor $(\delta(q, b), \omega)$ of some $(q, \omega) \in
C_j$ with $q \neq q_0$, or the fresh origin. In the first case take the realization $x = u\, W[0 \,..\, j)\, v$ of
$(q, \omega)$. Its covering factor reaches $q$ at the window's end. Replace its remainder $v$ by $b v'$, with $v'$ a
shortest completion from $\delta(q, b)$. That keeps every earlier factor and the origin. In the second case the seed
fired, so some $(q', \omega') \in C_j$ has $q'$ accepting. Take its realization and end its covering factor at the
window's end, which is legitimate since $q'$ accepts. Then begin a new factor $b v'$, with $v'$ a shortest completion
from $\delta(q_0, b)$. Its origin is $j$. The rename does not arise at $j \geq 1$: when $q_0$ is not re-entrant
nothing enters it, and when it is, the rename is off. The left context $u$ chosen at $j = 1$ is never changed. If $X$
is empty, no state is live and $C_k$ is empty, and no word of $X^*$ contains a non-empty window. So both sides are
empty.
\end{proof}

\begin{corollary}[What certification decides]
\label{cor:decides}
Let $W$ be a factor of a word of $X^*$. The model certifies $(W, o)$ exactly when, in every factorization of every
word of $X^*$ containing $W$, the factor covering the occurrence's last byte begins $o$ bytes into it.
\end{corollary}

The corollary explains the vacuous window of Section~\ref{sec:preliminaries}. Over $\{\code{0}, \code{00},
\code{01}\}$ the word \code{01001} factors as $\code{01} \mid \code{0} \mid \code{01}$. That factorization contains
\code{1001} at offset $1$, with the last factor beginning $2$ bytes into the window. Every factorization of a word
containing \code{1001} agrees. A \code{1} is only ever the second byte of \code{01}, so the window's two \code{1}'s
end factors, and the factor covering the last one is the \code{01} at offset $2$. The model therefore certifies
$(\code{1001}, 2)$, and rightly so for $X^*$. The scan of \code{01001} takes \code{01} and \code{00} and fails on
the last \code{1}. The word is in $X^*$ and not in $\mathcal{G}$. Non-emptiness of the cloud is an occurrence
theorem for $X^*$, not for $\mathcal{G}$. That is the whole gap between the model's certificates and the witnessed
ones.

\begin{corollary}[Token kinds are immaterial]
\label{cor:kinds}
Two automata with the same token language $X$ give, on every window, clouds with the same set of origins. Hence they
give the same certificates, the same origins and the same shortest certified length. Token kinds may be erased, the
positive-width equivalent taken, and the Boolean recognizer of $X$ minimized before any search of this paper.
Re-entrancy is recomputed on the result.
\end{corollary}

\begin{proof}
The set of origins of $F(W)$ depends only on $X$ and $W$. Certification, the origin and the shortest length are read
from that set.
\end{proof}

The reduction can be large. Give every binary word of length $k$ its own token kind. The kind-preserving minimal
automaton keeps every prefix apart and has $2^{k+1} - 1$ live states. The Boolean recognizer of $\{\code{a},
\code{b}\}^k$ has $k + 1$. Both give the same certificates. The evaluation runs on the lexer's own automata, as
shipped. The corollary is a reduction available to the searches, not one they take.

\subsection{Where the model is exact}
\label{sec:exact}

By Corollary~\ref{cor:decides} the model quantifies over factorizations, where Definition~\ref{def:window}
quantifies over scans. The two coincide exactly when factorizations and scans are the same thing.

\begin{corollary}[Exactness]
\label{cor:exact}
Suppose every word of $X^*$ has exactly one factorization and $\mathcal{G} = X^*$. Then for every window the origins
of the cloud are exactly the covering origins over the completely tokenizable occurrences of $W$, with $\before$
standing for every origin outside the window. The model certifies every witnessed certificate, and refuses every
occurring window that is not one. Every prefix code satisfies the hypotheses. So do $\{\code{a}, \code{ab}\}$ and
$\code{ab}^*$, which are not prefix codes.
\end{corollary}

\begin{proof}
Every word of $X^*$ is completely tokenizable, and its scan is a factorization, hence its only one. So the
factorizations of words of $X^*$ containing $W$ are exactly the scans of the completely tokenizable inputs
containing $W$. The origins of Theorem~\ref{thm:factorization} are then the covering origins. An occurring window
with one covering origin at every occurrence is certified by the model. One with two is refused, since both appear
in $F(W)$. A prefix code admits at most one codeword at any position. So a word of $X^*$ has one factorization, and
the scan computes it by taking the one codeword that begins. Over $\{\code{a}, \code{ab}\}$ every word of $X^*$ is a
sequence of \code{a}'s, each followed by at most one \code{b}. Its factors begin at the \code{a}'s, so the
factorization is unique. The scan computes it, taking \code{ab} where a \code{b} follows and \code{a} otherwise.
Over $\code{ab}^*$ the same holds with runs of \code{b}.
\end{proof}

Unique factorization alone is not enough; the scan must also succeed on all of $X^*$. Over $\{\code{a}, \code{ab},
\code{bb}\}$ every word of $X^*$ has one factorization. The reversed set $\{\code{a}, \code{ba}, \code{bb}\}$ is a
prefix code, so a word is factored uniquely from its right end. The window \code{ab} is semantically certified at
origin $0$ and witnessed by the input \code{ab}. An \code{a} is only ever a token's first byte, and the scan prefers
\code{ab} to \code{a} where a \code{b} follows. Yet \code{abb} factors as $\code{a} \mid \code{bb}$. That covers the
window's last byte with a factor beginning at offset $1$, so the model refuses the window. The scan of \code{abb}
takes \code{ab} and fails on the final \code{b}. So $\code{abb} \in X^* \setminus \mathcal{G}$, and that is exactly
the hypothesis that fails. The paper's witness one, $\{\code{a}, \code{ab}, \code{b}\}$, fails the other
hypothesis, since \code{ab} has two factorizations.

The hypotheses of Corollary~\ref{cor:exact} are decidable in polynomial time. Write $\mathrm{Pref}_+(X^*)$ for the
non-empty prefixes of words of $X^*$.

\begin{theorem}[The exact class]
\label{thm:agreement}
The following are equivalent for a token language $X$.
\begin{enumerate}
\item Every factorization of every word of $X^*$ is the scan of that word, which succeeds.
\item Every word of $X^*$ has exactly one factorization and $\mathcal{G} = X^*$.
\item $X \cap X\, \mathrm{Pref}_+(X^*) = \varnothing$: no word of $X$ is a word of $X$ followed by a non-empty prefix
of a word of $X^*$.
\end{enumerate}
Condition (3) is decided on the live automaton in $O(256 \cdot n^2)$ time. When it fails, the decision returns $u
\in X$, a non-empty $r$ with $u r \in X$, and $z$ with $r z \in X^*$, of total length at most $n^2 + 2n - 2$. The
word $u r z$ is then in $X^*$, and its factorization beginning with $u$ is not its scan.
\end{theorem}

\begin{proof}
(1) implies (2): a word of $X^*$ has a factorization, so its scan succeeds, and every factorization is that scan, so
there is one. (2) implies (1): the scan of a word of $X^*$ succeeds and is a factorization, hence the only one.
Suppose (3) fails: $u \in X$, $r$ is non-empty, $u r \in X$ and $r z \in X^*$. Then $u r z \in X^*$ has a
factorization beginning with $u$. Its scan, if it succeeds, begins with a token of length at least $|u r| > |u|$.
Either the scan fails or it is a different factorization, so (1) fails. Suppose (3) holds and take any factorization
$x = x_1 \cdots x_m$. If $m = 0$ the scan of the empty word succeeds. Otherwise $x_1 \in X$ is a candidate first
token. A longer accepted prefix $x_1 r$ of $x$ would have $r$ a non-empty prefix of $x_2 \cdots x_m \in X^*$, against
(3). So the scan commits $x_1$ and restarts. By induction on $m$ it commits $x_2, \ldots, x_m$ and succeeds.

For the decision, $\mathrm{Pref}_+(X^*)$ is recognized by a nondeterministic automaton $P$ over the live states. Start
at $q_0$ and follow $\delta$. Before any byte read from an accepting state, optionally reset to $q_0$. Every state
reached after at least one byte accepts. A path of $P$ spells a sequence of accepted factors followed by a prefix of a
live run, which completes to a factor; so the word is in $\mathrm{Pref}_+(X^*)$. Conversely a non-empty prefix of $x_1
\cdots x_m$ is spelled by resetting at each factor's end. Now run one breadth-first search, started at every pair $(q,
q_0)$ with $q$ a reachable accepting state, over pairs $(p, p')$. Here $p$ continues $\delta$ from $q$ and $p'$ runs $P$
from $q_0$, both over the same bytes, and a pair reached after at least one byte is marked as started. A started pair
with $p$ accepting is an obstruction. Let $u$ be a shortest word reaching the $q$ its search path began at, of length at
most $n - 1$ since $q_0$ does not accept. Let $r$ be the bytes read, and $z$ a shortest accepting completion from $p'$,
of length at most $n - 1$. Then $u \in X$, $u r \in X$ and $r z \in X^*$. Conversely an obstruction $u$, $r$ drives the
search from $\delta^*(q_0, u)$ along $r$ to such a pair. There are at most $n^2$ started pairs, and a shortest path to
one repeats none, so $|r| \leq n^2$. The search visits each pair's at most two successors per byte once.
\end{proof}

The figure program runs the decision on the examples above. $\{\code{a}, \code{ab}\}$ and $\code{ab}^*$ are in the
class. $\{\code{a}, \code{ab}, \code{bb}\}$ is out with the obstruction \code{abb}, and $\{\code{a}, \code{ab},
\code{b}\}$ with \code{ab}. $\{\code{0}, \code{00}, \code{01}\}$ is out with \code{00}, whose factorization
$\code{0} \mid \code{0}$ the scan does not take. The conventional C-like row is out with \code{//} followed by a
tab: a comment and a prefix of a whitespace run, which the scan reads as one comment. The class is a sufficient
condition for the model to be exact, not a necessary one. Exactness compares unanimous origins, a weaker observable
than equality of every history.

\subsection{A shortest witnessed window}
\label{sec:witnessed}

The evaluation separates model certificates from witnessed ones by a bounded search for an occurrence. A certificate
it cannot witness is reported as unresolved. The separation can be decided exactly. A shortest witnessed window is
found, with its occurrence, by running the chosen-origin automaton of Section~\ref{sec:chosen} in step with a
recognizer of the factors of $\mathcal{G}$. The recognizer is the companion's armed-run
automaton~\cite{nidhogg2026splitting}, which we use as a black box. An input is
completely tokenizable exactly when it segments into accepted words such that no committed word extends, within the
input, to a longer accepted word from the same start. The automaton guesses the boundaries, runs each committed
word's state onward as an armed run, and rejects if an armed run accepts again. Its states are a current token state
or a fresh boundary, together with the set of armed states, at most $(n + 1)\, 2^n$ of them. It recognizes exactly
$\mathcal{G}$.

Let $H$ be that automaton, trimmed to the states reachable from its start and co-accessible to its accepting states.
Let $\mathrm{Fact}(H)$ be $H$ with every retained state made initial and accepting.

\begin{lemma}[Factors]
\label{lem:factors}
$\mathrm{Fact}(H)$ accepts exactly the factors of completely tokenizable inputs.
\end{lemma}

\begin{proof}
Take a path of $H$ between retained states. The trim gives a stem from the start to its first state and a suffix
from its last state to an accepting state. The three together are an accepting run of $H$, so the path's label is a
factor of a word of $\mathcal{G}$. Conversely a factor of a word of $\mathcal{G}$ is the label of the middle of an
accepting run of $H$, whose states are all retained.
\end{proof}

\begin{theorem}[Witnessed synthesis]
\label{thm:witnessed}
The product of the chosen-origin automaton with $\mathrm{Fact}(H)$ accepts exactly the witnessed model certificates.
Breadth-first search over it decides whether one exists. It returns a shortest one with its origin and a completely
tokenizable input containing it, and its exhaustion proves that no witnessed model certificate exists at any length.
The product has at most $K (n + 1)\, 2^n = O(n^2 4^n)$ states. Existence of a witnessed model certificate is in
PSPACE.
\end{theorem}

\begin{proof}
The product accepts the intersection of the two languages: the model-certified windows by Theorem~\ref{thm:chosen},
and the factors of $\mathcal{G}$ by Lemma~\ref{lem:factors}. That intersection is the set of witnessed model
certificates, by Definition~\ref{def:witnessed}. Every edge reads one byte, so breadth-first search from the initial
pairs finds a shortest accepted window, and the chosen-origin component gives its origin. For the occurrence, take
the $H$ states at the ends of the accepting path. Attach the stem $u$ and the suffix $v$ the trim guarantees. Then
$u W v$ is accepted by $H$, so it is completely tokenizable, and it contains $W$ at offset $|u|$. By
Theorem~\ref{thm:soundness} the token covering that occurrence's last byte begins at the certified origin. For the
space bound, a product state is written in $O(n)$ bits. A path of length below the product's size is guessed with a
counter of $O(n)$ bits. The initial and final $H$ states are checked reachable and co-accessible by implicit
searches in $H$ within the same space. So the problem is in nondeterministic polynomial space.
\end{proof}

Exhaustion here means no witnessed \emph{model} certificate. It still says nothing about a semantic certificate the
model refuses. The returned input is a witness of occurrence, not a shortest such input, since the search minimizes
the window. The four model-positive grammars of the random sweep that the bounded witness search leaves unresolved
are a question this search settles either way. This paper does not rerun the sweep under it. With
Lemma~\ref{lem:reduction-forward}, whose window is witnessed, the witnessed half of Theorem~\ref{thm:pspace}
follows.

\subsection{Whether the model misses a certificate}
\label{sec:completeness}

Section~\ref{sec:strictness} shows the model refusing a witnessed certificate over two small token sets. Whether a
given token set has any such refusal is decidable. Call the model \emph{globally complete} for a token set when it
certifies every witnessed certificate of Definition~\ref{def:witnessed}. Completeness holds vacuously when there is
none.

Two marked graphs carry the decision. In both, an edge is labelled by a byte and marked when a token begins just
before that byte. The first is the trimmed $H$, whose paths are the greedy histories. The second, $\Phi$, has the live
states and a fresh boundary $\bot$. It has an unmarked edge $p \xrightarrow{b} \delta(p, b)$ for every live
transition, a marked edge $p \xrightarrow{b} \delta(q_0, b)$ whenever $p$ accepts, and marked edges $\bot
\xrightarrow{b} \delta(q_0, b)$. Its initial state is $\bot$ and its final states are the accepting live states, and
it is trimmed likewise. Its accepting paths are exactly the factorizations of words of $X^*$. By the argument of
Lemma~\ref{lem:factors}, its paths between retained states are exactly the factors of words of $X^*$, with the
factor starts they contain marked. Write $N$ and $f$ for the two graphs' state counts, with $N \leq (n + 1)\, 2^n$
and $f \leq n + 1$. For a graph $J$ and a candidate $(W, o)$, a path spelling $W$ \emph{has origin $o$} when its edge
at byte $o$ is marked and no later edge is. Let $\mathrm{Occ}_J$ be the set of candidates some path of $J$ spells.
Let $\mathrm{Bad}_J$ be the set of candidates some path of $J$ spells with a different origin.

\begin{theorem}[Global completeness]
\label{thm:complete}
The model misses $(W, o)$, a witnessed certificate it refuses, exactly when $(W, o) \in \mathrm{Occ}_H \cap
\mathrm{Bad}_\Phi \setminus \mathrm{Bad}_H$. Whether the model is globally complete is decidable. When it is not, a
shortest missed window is returned with its origin, a completely tokenizable input containing it, and a word of $X^*$
with a factorization that covers the occurrence's last byte from another origin. The decision runs over an observer
of at most $2^{N + f} + 4^{N + f}$ states. So it is in EXPSPACE in $n$, and takes doubly exponential time.
\end{theorem}

\begin{proof}
A witnessed certificate $(W, o)$ is a candidate in $\mathrm{Occ}_H$, since $W$ occurs in some completely tokenizable
input. It is outside $\mathrm{Bad}_H$, since every greedy occurrence has origin $o$. The model refuses it exactly
when $C_k$ carries an origin other than $o$. Indeed the pair of the greedy occurrence is in $C_k$ by
Lemma~\ref{lem:representation} and carries $o$, so the cloud is non-empty and certification fails only by a second
origin. By Theorem~\ref{thm:factorization} a second origin is a factorization of a word of $X^*$ covering the
occurrence's last byte from elsewhere, which is membership in $\mathrm{Bad}_\Phi$. So the three conditions are exactly
the missed certificates.

The observer reads $W$ with its byte $o$ marked. For each graph $J$ it holds the set of ends of the paths spelling
the prefix read so far. After the mark that set is split in two: the \emph{right} set $R_J$ of ends whose path has
origin $o$ so far, and the \emph{wrong} set $B_J$ of the rest. Before the mark it holds one set per graph, initially
every retained state. At the marked byte, an edge from the set goes to $R_J$ when marked and to $B_J$ otherwise.
After it, $R_J'$ is the set of ends of unmarked edges from $R_J$. $B_J'$ is the set of ends of all edges from $B_J$,
together with the ends of marked edges from $R_J$. The two sets may overlap, since different paths reaching one
state may differ in origin, and they are held apart. At the window's end, $(W, o)$ is missed exactly when $R_H \neq
\varnothing$, $B_H = \varnothing$ and $B_\Phi \neq \varnothing$, which restates the three conditions. Given the position
of the mark the observer is deterministic, and the mark is one extra choice per window. So breadth-first search over
observer states, marking at each depth or not, reaches a missed window of minimum length or exhausts the finitely
many states. From a missed window, a path through $H$ into $R_H$ and a path through $\Phi$ into $B_\Phi$ are read back.
Stems and suffixes complete them into the two inputs claimed. A state of the observer is one or two subsets of each
graph, hence the bound. It is written in $O(N + f)$ bits, which is $O(n 2^n)$. A nondeterministic search with a
counter runs in that space, and the explicit search runs in time polynomial in the observer's size.
\end{proof}

Over $\{\code{a}, \code{ab}, \code{b}\}$ the shortest missed window is \code{ab} at origin $0$. The completely
tokenizable input is \code{ab}, and the factorization is $\code{a} \mid \code{b}$ of the same word. Over $\{\code{a},
\code{ab}, \code{bb}\}$ it is the same window and input, with the factorization $\code{a} \mid \code{bb}$ of
\code{abb}. Over $\{\code{a}\}$ and over $\{\code{a}, \code{ab}\}$ the model is globally complete, the latter by
Corollary~\ref{cor:exact}. The bound is an upper bound only. Whether a polynomial criterion for global completeness
exists is open; Theorem~\ref{thm:agreement} decides a sufficient condition in polynomial time. Completeness is a
property of the token set, settled once from its tables. Where it holds, every negative of the model is a negative
of Definition~\ref{def:witnessed}, and the three-layer reading of the results collapses to two.

\section{From a boundary to a parallel cut}
\label{sec:cut}

At every occurrence at offset $t$ in a completely tokenizable input, a certified window $(W, o)$ yields a true boundary
$t + o$ of the final segmentation. Turning a boundary into a parallel cut is the predecessor's territory. The
predecessor's prefix-stability result is what licenses a worker to scan from a known boundary and agree with the serial
stream around the cut~\cite{nidhogg2026splitpoints}. The division of labour is exact: this paper establishes that $t +
o$ is a boundary; the predecessor establishes what a scan starting at a boundary preserves. The v1.3.3 artifact
evaluated here plans with single-byte certificates only. Release 1.4.0 added the window decision as a public contract,
which the sweep at the later commit cross-checks against its model. The same release added a window-planning sibling,
which lies outside this paper's evaluation and is not a claim of this paper.

\section{Boundary profiles}
\label{sec:profile}

Section~\ref{sec:preliminaries} kept three guarantees apart, and this paper certifies the strongest. This section asks
the weakest, a fixed safe boundary, of every gap of a window at once. The answer places the certificate among the
anchors a window can soundly supply. The property is the one release 2.1.0 of munch decides exactly within its search
cap.

\begin{definition}[Boundary profile]
\label{def:profile}
Let $W = w_0 \cdots w_{k-1}$ with $k \geq 1$. A \emph{gap} of $W$ is an index $g \in \{0, \ldots, k\}$: the gap before
byte $w_g$ when $g < k$, and the gap after the final byte when $g = k$. Take an occurrence of $W$ at offset $t$ in a
completely tokenizable input $x$. The gap $g$ \emph{is a boundary} there when a token of the maximal-munch tokenization
of $x$ begins at $t + g$ or $t + g = |x|$, and a token \emph{crosses} it otherwise. Over every occurrence of $W$ in
every completely tokenizable input, the gap is \emph{forced} when it is a boundary at each of them, \emph{crossed} when
it is a boundary at none, and \emph{open} when both happen. The \emph{boundary profile} of $W$ assigns every gap its
verdict.
\end{definition}

A window no completely tokenizable input contains is forced and crossed at every gap, both vacuously, exactly as
Definition~\ref{def:window} holds vacuously there. For an occurring window each gap has one verdict. The library names
the three verdicts \code{must}, \code{never} and \code{may}, and a window occurring nowhere \code{absent}.

\begin{remark}
\label{rem:profile}
Three earlier notions are profile patterns by definition. The pair $(W, o)$ is a certified split window exactly when gap
$o$ is forced and gaps $o + 1, \ldots, k - 1$ are crossed. At an occurrence, the token covering byte $t + k - 1$ begins
at $t + o$ exactly when a token begins at $t + o$ and none begins at $t + o + 1, \ldots, t + k - 1$. Both conditions are
quantified over the same occurrences, so the conjunction splits gap by gap. A certified split point of the
predecessor~\cite{nidhogg2026splitpoints} is the case $k = 1$ with gap $0$ forced. The fixed safe boundary of
Section~\ref{sec:preliminaries} is a forced gap.
\end{remark}

\paragraph{Local maximality.} The weakest guarantee is the one a cut consumes. Fix a width $H \geq 1$. A \emph{cut rule
of width $H$} is a set $R$ of pairs $(W, g)$ with $1 \leq |W| \leq H$ and $g$ a gap of $W$. On an input $x$ the rule
cuts at every position $t + g$ such that $(W, g) \in R$ and $W$ occurs at $t$. Each piece between consecutive cuts is
scanned from $q_0$ on its own, and the token streams are concatenated with nothing reconciled. The rule is \emph{sound}
when on every completely tokenizable input the concatenation is the input's maximal-munch tokenization. The rule reads
the window and nothing else: no state crosses a cut, and no piece is rescanned.

\begin{proposition}[Local maximality]
\label{prop:maximal}
A cut rule is sound if and only if every gap it cuts at is forced.
\end{proposition}

\begin{proof}
Suppose every $(W, g) \in R$ has gap $g$ forced, and fix a completely tokenizable $x$. Every cut falls at a forced gap
of an occurrence, so at a boundary of the tokenization of $x$ or at $|x|$. A cut at a boundary is exact: the pieces
between boundaries, scanned from $q_0$ and concatenated, reproduce the tokenization. That is the companion's lemma that
cuts at boundaries are exact~\cite{nidhogg2026splitting}, stated there for literal token sets and carried to regular
ones by its corollary on splitting at the regular generality. Section~\ref{sec:cut} composes with the same step.
Conversely, suppose some $(W, g) \in R$ has gap $g$ not forced. Some completely tokenizable $x$ then contains $W$ at an
offset $t$ where a token of the tokenization crosses $t + g$, beginning before it and ending after it. The rule cuts $x$
at $t + g$. The piece ending there emits no token reaching past $t + g$, and the piece beginning there emits none
beginning before it. The crossing token is emitted by neither, so the concatenation differs from the tokenization.
\end{proof}

The certificate is the case of one forced gap followed by crossed gaps to the window's end, and it supplies the anchor
$t + o$. The proposition says what else a window supplies. A sound rule of width $H$ cuts at no position the forced gaps
of windows of width at most $H$ fail to name, and it may cut at every position they name. The profile at width $H$ is
therefore the whole anchor supply of that width. The certificate names one member of it per window. The gap $k$ after
the window is the member no origin can name: a forced gap $k$ of $W$ is the certificate $(Wc, k)$ for every byte $c$
that can follow, one byte later.

\paragraph{Decidability.} Gap $g$ is not forced exactly when some completely tokenizable input has an occurrence of $W$
that a token crosses at $g$. The gap is not crossed exactly when some has an occurrence cut there. Each is a question of
the form the companion's offline verifiability theorem decides~\cite{nidhogg2026splitting}. The boundary markings of
completely tokenizable inputs form a regular language, recognized by its armed-run automaton. Its intersection with the
markings holding an occurrence of $W$ with, or without, a mark at gap $g$ has decidable emptiness with a shortest
witness. The profile is therefore decidable for every regular token set, and nothing is proved fresh here. Release 2.1.0
of munch decides both questions per gap, as \code{boundary\_counterexample()} and \code{crossing\_counterexample()}.
Both run the one search its exact window decision \code{window\_counterexample()} runs: the input is built byte by byte,
breadth first, with its token boundaries guessed and every guess held to maximal munch. Each returns the shortest
refuting input or, once it has exhausted its states under a cap on the states held at once, the proof that none exists.
A search the cap stops settles nothing, and the gap is reported undetermined. \code{boundary\_profile()} asks both
questions at every gap and reads the verdicts off. The gap is \code{must} when only the first search exhausts without a
witness, \code{never} when only the second does, and \code{may} when both find one. Whether the window occurs at all is
decided once per window by \code{window\_occurrence()} under the same cap, and \code{absent} is given at every gap or at
none.

\paragraph{Anchor propagation.} Section~\ref{sec:preliminaries} noted that the weakest guarantee is already usable: a
worker cuts at the known boundary and rescans the window's suffix. Read as a window, that rescan forces further gaps.
The statement needs a longest token. Write $L$ for the maximum token length of a literal token set, a finite set of byte
strings, or of a regular token set whose token languages are all finite. A regular token set with an infinite token
language has no $L$, and the lemma is not stated for it. Two facts precede it. First, extension is monotone. A forced or
crossed gap of $W$ stays forced or crossed, shifted, in every window that contains $W$, since every occurrence of the
larger window contains an occurrence of $W$ at the shifted offset. Second, the scan restarts in $q_0$ at every boundary,
which is the companion's memorylessness lemma~\cite{nidhogg2026splitting}. From a boundary $p$ of a completely
tokenizable $x$ it therefore commits the longest token that is a prefix of $x[p \,..)$.

\begin{lemma}[Anchor propagation]
\label{lem:propagation}
Let the token set have maximum token length $L$, let $W$ occur in some completely tokenizable input, and let gap $g$ of
$W$ be forced. Define offsets $g = c_0 < c_1 < \cdots < c_m$ by $c_{i+1} = c_i + |\tau_i|$ for as long as $c_i + L \leq
k$, where $\tau_i$ is the longest token that is a prefix of $W[c_i \,..\, k)$. Then every $c_i$ is a forced gap of $W$,
and every gap strictly between $c_i$ and $c_{i+1}$ is crossed.
\end{lemma}

\begin{proof}
Take an occurrence of $W$ at offset $t$ in a completely tokenizable $x$, and suppose $t + c_i$ is a boundary of the
tokenization of $x$, as it is for $i = 0$ by hypothesis. The scan commits at $t + c_i$ the longest token prefixing $x[t
+ c_i \,..)$. Every token has length at most $L$, so that token is determined by $x[t + c_i \,..\, t + c_i + L)$, which
is $W[c_i \,..\, c_i + L)$ when $c_i + L \leq k$. The committed token is therefore $\tau_i$, which exists because the
occurrence is tokenized, and it is the same at every occurrence. Its end $t + c_{i+1}$ is a boundary, and no boundary
lies strictly inside it. Induction on $i$ gives both claims.
\end{proof}

\begin{corollary}[Resolving width]
\label{cor:width}
Let the token set have maximum token length $L$, let $x$ be completely tokenizable, and let $b$ be a boundary of its
tokenization with $b + L \leq |x|$. Let $a < b$ be a boundary that some window $W_a$ of width $h_a$ forces: $W_a$ occurs
at $q$, and $a = q + g_a$ for a forced gap $g_a$. Then the window $x[q \,..\, e)$ with $e = \max(b + L,\, q + h_a)$
forces $b$ at its gap $b - q$. A window of width at most $(b - a) + h_a + L$ therefore forces $b$.
\end{corollary}

\begin{proof}
The window contains $W_a$ at offset $0$, so its gap $g_a$ is forced by monotonicity, and it occurs in $x$. Let $a = p_0
< p_1 < \cdots < p_m = b$ be the boundaries of $x$ from $a$ to $b$. Each is the end of the token the scan commits at the
one before, and each $p_i$ with $i < m$ satisfies $p_i + L < b + L \leq e$. In the window's offsets these are the $c_i$
of Lemma~\ref{lem:propagation}, each with $c_i + L \leq e - q$, so $c_m = b - q$ is forced. The width is $e - q =
\max((b - a) + g_a + L,\, h_a) \leq (b - a) + h_a + L$.
\end{proof}

In words, a boundary within distance $d$ of an anchor is forced by a window of width about $d$ that reaches back to the
anchor. That is the guarantee of Section~\ref{sec:preliminaries}, cut at the known boundary and rescan the suffix, read
as a window: the window carries the anchor, and the rescan is the lemma's induction. The bound is not the only route to
a forced gap. The measurements below show windows forcing a boundary with no anchor in reach, from local structure the
lemma does not describe.

\paragraph{Reach and its limits.} Fix an input $x$ and a width $H$. The \emph{anchors of $x$ at width $H$} are the
positions $t + g$ with $0 < t + g < |x|$ such that some window $x[t \,..\, t + h)$ with $h \leq H$ has gap $g$ forced.
The \emph{reach} $R_H(x)$ is the length of the longest stretch of $x$ between consecutive anchors, with $0$ and $|x|$
counting as anchors. The \emph{all-width reach} $R_\infty(x)$ admits every finite width, and the \emph{floor} $F(x)$ is
the length of the longest token of the tokenization of $x$. Then
\[ R_H(x) \;\geq\; R_\infty(x) \;\geq\; F(x). \]
A wider budget adds anchors and removes none. Every anchor is a boundary, so none lies strictly inside the longest
token, and some stretch contains that token whole. The floor is shared by every method that cuts only at boundaries.
Below it only a cut inside a token, resumed from scanner state, can go. Two facts locate the middle term.

% The bound's check: explorations/cut-contexts/logs/reach-bound-check.txt, written by reach_bound_check.py there: L1
% (the example: boundaries 0 2 4 5 8 9, anchors 5 8, R_inf 5, F 3), L4 (82336 inputs checked), L5 (12526 above the
% floor), L7 (0 bound violations), L8 (0 corollary violations).
(a) With a maximum token length $L$, Corollary~\ref{cor:width} makes every boundary $b$ past an anchor $a$, with $b + L
\leq |x|$, an anchor at some finite width. Past its first anchor and up to $L$ bytes before its end, the all-width
stretches of such an input are single tokens. The stretch before the first anchor need not be. Over $\{\code{a},
\code{aa}, \code{bbb}\}$ the input \code{aaaaabbba} has boundaries $0, 2, 4, 5, 8, 9$ and, at every width, the anchors
$5$ and $8$ alone, since a longer run of \code{a} may precede any occurrence. The all-width reach is $5$ against a floor
of $3$. What the corollary gives is the bound
\[ R_\infty(x) \;\leq\; \max(a_1,\, F(x) + L - 1), \]
with $a_1$ the first anchor of $x$, taken as $|x|$ when there is none. The head stretch has length $a_1$. Every other
stretch begins at an anchor, where the scan commits a token ending at a boundary $b$. Either $b$ is an anchor and the
stretch is that token, or $b + L > |x|$ by the corollary and the stretch ends within $L - 1$ bytes past $b$. A
brute-force check over every literal set of one to three tokens of length at most $3$ over $\{\code{a}, \code{b}\}$, and
every completely tokenizable input of length at most $10$, found no violation in 82{,}336 inputs, 12{,}526 of them above
the floor. Past the head, the middle term measures how far apart the anchors of the input lie, not an ambiguity wider
windows cannot resolve.

(b) No uniform radius exists for the JSON row. For $H \geq 1$ let $v_H$ be the prefix of length $H$ of the periodic byte
string $\code{,1,1,1}\cdots$. The input $v_H$ itself tokenizes into one-byte tokens, a comma or the number \code{1}
each, so every gap of that occurrence is a boundary. The input $\code{"}v_H\code{"}$ is one string token, since neither
byte is a quote, a backslash or a control byte, so a token crosses every gap of that occurrence. Both inputs are
completely tokenizable, so every gap of $v_H$ is open, at every width $H$. The two occurrences differ in whether an
unmatched quote lies to their left, which no bounded window sees. The conventional and split-friendly rows hold the same
family between the quotes of a string literal: outside a string a comma is a punctuation token and \code{1} a number,
and the string form excludes only a quote and a newline. The block-comment rows hold it between \code{/*} and \code{*/},
which the family never closes, having no \code{*}. The bare row has neither strings nor comments, so the family says
nothing there. In the terms of constrained coding, read the token set as a tagged graph: the companion's armed-run
automaton, each edge labelled by its byte and tagged by its boundary bit. A uniform radius $r$ is the tagged analogue of
the $(m, a)$-sliding-block decodability of Marcus, Roth and Siegel~\cite[Section~4.3]{marcus2001constrained}, defined
there for tagged encoders, which the armed-run automaton is not. Any two paths generating the same word agree on the tag
of their central edge, and the radius maps to $(m, a) = (r, r - 1)$ when that edge is the byte after the gap: $r$ bytes
precede it and $r - 1$ follow it. The family above is its failure for every look-behind and look-ahead: the quoted and
the bare reading are two paths over the same bytes with opposite tags throughout. Deciding that ceiling for a token set
is an analysis of the cycles of the pair graph of the companion's verifier, two copies run over equal bytes, which this
paper does not carry out.

\paragraph{Measurements.} The campaign decided every distinct substring of width $1$ to $H$ of each corpus, nothing
sampled, both ways: the certificate exactly, by \code{window\_counterexample()} at every origin, and the profile by
\code{boundary\_profile()}. Every search ran under the library's default cap of $2^{20}$ states, and none reached it. An
anchor is a position strictly inside the corpus that a certified origin, or a forced gap, of an occurring window names,
counted once. Every anchor was checked against the sequential scan to be a token boundary, and every crossed gap was
checked at every occurrence. The record is the campaign note and its outputs in the author's research repository, which
is private: \path{notes/cut-contexts-campaign-2026-09-28.md} and \path{explorations/cut-contexts/}. The figures are
exploratory until an archived record exists, and the probe of Section~\ref{sec:evaluation} asserts none of them.

% Absent verdicts: check_profile_absent = 0 at clike-*-H4.summary.txt L39, json-twitter-H4.summary.txt L39,
% local-384-H4.summary.txt L39, local-384-H3.summary.txt L37, gpt2-4k-H3.summary.txt L37 and every H5 summary L41. The
% JSON lexer: logs/digests.txt L15-16 (libs/json/src/lexer.cpp of hopper at HEAD caff92b), identical to v1.0.0's.
The library is munch, commit \code{df818b6a224221cf964095dd8e2c20735f506701}, one commit before the absence verdict
moved to one decision per window. The move changed nothing here, since no run reported an absent gap. The JSON row runs
over \code{twitter.json} of simdjson v3.10.1 under the JSON lexer of hopper, the author's parser library~\cite{hopper}:
\path{libs/json/src/lexer.cpp} at commit \code{caff92b}, unchanged since release v1.0.0, compiled into the probe as it
stands. Its shortest input segmented differently from the recovery companion's JSON row, found by the library's
segmentation decision, is a string literal holding an ill-formed UTF-8 byte. The C-like rows run over the 512 KiB
generated corpora of the recovery companion's campaign~\cite{nidhogg2026panicmode}, each under a token set the same
decision proves equivalent to the corpus's generator. Two literal vocabularies of the splitting companion's supply study
run over a 16 KiB slice of \code{twitter.json}: a byte-pair vocabulary of 384 merges trained on that corpus, 640 tokens,
and the first 4{,}096 merges of the GPT-2 vocabulary, 4{,}352 tokens. The width $H$ is $4$ for JSON and C-like and $3$
for the vocabularies. Table~\ref{tab:profile} gives the anchor density and the longest anchor-free stretch of both
detectors.

% Every figure of the table is printed in explorations/cut-contexts/out/<row>.summary.txt of the research repository at
% the line named beside its row (exact per KiB, must per KiB, exact longest free, must longest free), and the gain in
% logs/headline-H.txt at the line named, computed by headline.py from the same summary lines.
\begin{table}
\centering
\caption{Anchors from the certificate and from forced gaps at the campaign width $H$: distinct positions per KiB of
corpus, the gain of the profile over the certificate, and the longest anchor-free stretch in bytes under each. The
block-comment row is the campaign's token set proved equivalent to its corpus's generator; munch's own block-comment
grammar gives the same figures on that corpus.}
\label{tab:profile}
\setlength{\tabcolsep}{4pt}
\begin{tabular}{lcrrrrr}
\toprule
 & & \multicolumn{2}{c}{Anchors per KiB} & & \multicolumn{2}{c}{Longest stretch} \\
\cmidrule(lr){3-4} \cmidrule(lr){6-7}
Token set & $H$ & certificate & forced gaps & Gain & certificate & forced gaps \\
\midrule
% json-twitter-H4.summary.txt L151, L158, L156, L163; gain logs/headline-H.txt L3
JSON, \code{twitter.json} & 4 & 94.64 & 123.73 & 30.7\% & 465 & 465 \\
% clike-conventional-H4.summary.txt L151, L158, L156, L163; gain logs/headline-H.txt L4
C-like, conventional & 4 & 32.01 & 32.01 & 0 & 104 & 104 \\
% clike-split-friendly-H4.summary.txt L151, L158, L156, L163; gain logs/headline-H.txt L5
C-like, split-friendly & 4 & 47.33 & 47.33 & 0 & 104 & 104 \\
% clike-block-comments-alone-H4.summary.txt L151, L158, L156, L163; gain logs/headline-H.txt L7; the same figures at
% the same lines of clike-block-comments-H4.summary.txt, gain logs/headline-H.txt L6
C-like, block comments & 4 & 6.71 & 6.71 & 0 & 965 & 965 \\
% clike-bare-H4.summary.txt L151, L158, L156, L163; gain logs/headline-H.txt L8
C-like, bare & 4 & 324.42 & 324.42 & 0 & 10 & 10 \\
% local-384-H3.summary.txt L116, L123, L121, L128; gain logs/headline-H.txt L9
Byte-pair, 384 merges & 3 & 230.69 & 240.31 & 4.2\% & 61 & 61 \\
% gpt2-4k-H3.summary.txt L116, L123, L121, L128; gain logs/headline-H.txt L10
GPT-2, 4{,}096 merges & 3 & 620.88 & 630.75 & 1.6\% & 13 & 13 \\
\bottomrule
\end{tabular}
\end{table}

% The C-like gain below the campaign width, every position via the end gap: clike-conventional-H4.summary.txt L72-74
% (width 1: 0), L105-107 (width 2: 2140, all via the end gap), L138-140 and L171-173 (widths 3 and 4: 0);
% clike-split-friendly-H4.summary.txt L72-74 (10506), L105-107 (2140), L138-140 and L171-173 (0);
% clike-bare-H4.summary.txt L72-74 (26499), L105-107, L138-140 and L171-173 (0);
% clike-block-comments-alone-H4.summary.txt L72-74, L105-107, L138-140 and L171-173 (0 throughout). JSON's split:
% json-twitter-H4.summary.txt L171-173 (17940 profile-only positions, 16517 via an inner gap). The example window is
% derived from the token set in the text; the campaign's per-window output, which is not committed, agrees (profile
% mNMMm, exact rrrC). The 465-byte stretch is one string: rw-json-twitter.summary.txt L14-17 and L27 (no stretch longer
% than the floor). local-384 at width 4: local-384-H4.summary.txt L156 (exact 61) and L163 (must 40), the floor
% rw-local-384.summary.txt L14 (40). The model against the exact certificate: local-384-H3.summary.txt L109 (224.88) and
% L116 (230.69); gpt2-4k-H3.summary.txt L109 (615.38) and L116 (620.88).
The gain has two sources. A forced gap $k$ is the certificate of the window's one-byte extensions, so it is what the
next width buys the certificate. The C-like rows gain only below the campaign width, and only this way: the conventional
row gains 2{,}140 positions at width $2$, the split-friendly row 10{,}506 at width $1$ and 2{,}140 at width $2$, the
bare row 26{,}499 at width $1$, and the block-comment row none, every one from the gap after the window. The gain is
gone from width $3$ on the first two rows and from width $2$ on bare, and the profile shortens no C-like stretch at the
campaign width. The other source is several boundaries per window, which no certificate names: 16{,}517 of JSON's
17{,}940 profile-only positions come from a gap inside the window. One such window is two spaces, a closing bracket and
a newline. No string holds a raw newline, so the occurrence lies outside every string, the bracket is a token and the
newline begins a whitespace run: gaps $2$ and $3$ are forced, while the certificate names $3$ alone. The 465-byte
stretch on JSON is one string literal, the floor itself, which nothing at any width can shorten. On the vocabularies the
gain is 4.2 and 1.6 percent. At width $4$ the 384-merge vocabulary's longest stretch under forced gaps falls from 61 to
40 bytes, the floor, while under the certificate it stays at 61. The model of Section~\ref{sec:model} equals the exact
certificate on every JSON and C-like row. The model falls short only on the vocabularies, at 224.88 against 230.69
anchors per KiB on the 384-merge vocabulary and 615.38 against 620.88 on the GPT-2 prefix.

% Floors: rw-<row>.summary.txt L14 (json 465, clike-conventional 15, clike-split-friendly 15, clike-block-comments-alone
% 135, clike-bare 10, local-384 40, gpt2-4k 9); stretches longer than the floor: L27 (json 0, bare 0). Conventional's
% 104-byte stretch with 17 starts: rw-clike-conventional.summary.txt L30, L32; block comments' 965 with 283:
% rw-clike-block-comments-alone.summary.txt L30, L32. Resolution ladders: rw-clike-conventional.summary.txt L75-81 and
% rw-clike-split-friendly.summary.txt L75-81 (unresolved 0 at 128); rw-clike-block-comments-alone.summary.txt L81
% (33 unresolved at 128); rw-clike-block-comments-alone-W1024.summary.txt L84 (0 unresolved at 1024). The corpora:
% logs/digests.txt L9-10 (one digest for conventional and split-friendly). Reach-back, counted as reaching to within a
% byte of the anchor or past it: logs/rw-reachback.txt L1 (conventional 40 of 40, 5..95), L2 (split-friendly 40 of 40,
% 5..96), L4 (block comments 40 of 40, 1..912), L5 (local-384: 11 of 12 not reaching back, distances 3..45), L6
% (gpt2-4k 6 of 6 reaching back, 3..12). The two boundaries: logs/rw-clike-conventional.err L22 (b 187441, guided
% width 106) against rw-clike-conventional.summary.txt L28 (stretch begins at 187346, so 95 past), and
% logs/rw-clike-block-comments-alone-W1024.err L36 (b 217108, guided width 933) against
% rw-clike-block-comments-alone-W1024.summary.txt L28 (begins at 216196, so 912 past). The curves:
% rw-clike-conventional.summary.txt L92, L95, L98, L101, L104, L107 (104, 101, 91, 82, 44, 16) and
% rw-clike-block-comments-alone-W1024.summary.txt L96, L111, L120 (965, 857, 127). Vocabulary ladders:
% rw-local-384.summary.txt L75-77 (2, 7, 3 by widths 4, 8, 16) and rw-gpt2-4k.summary.txt L75, L77 (3 and 3).
A second run asked how wide a window must be to force the boundaries inside each corpus's largest anchor-free stretches.
Up to 40 boundaries per row, evenly spaced, were asked of \code{boundary\_counterexample()} over a ladder of widths $4$
to $128$ with three splits each, and of a witness-guided growth beside it, every answer checked against the scan. JSON
and bare are at the floor: no stretch is longer than the longest token, 465 bytes of one string on JSON and 10 on bare,
so nothing was searched. The C-like rows are not. The conventional row's longest stretch is 104 bytes holding 17 token
starts, where the longest token is 15 bytes, and the block-comment row's is 965 bytes holding 283, where the longest
token is 135 bytes. The conventional and split-friendly corpora are byte-identical, which is why both rows show 104.
Every searched boundary resolves: all 40 on the conventional and split-friendly rows by width 128, and on block comments
7 of 40 by 128 and all 40 by 1{,}024. Every resolving window on a C-like row reaches back to, or past, the stretch's
left anchor. That is 40 of 40 on the conventional row over distances of 5 to 95 bytes, 40 of 40 on split-friendly over 5
to 96, and 40 of 40 on block comments over 1 to 912. A boundary 95 bytes past its anchor resolves with a guided window
of 106 bytes, and one 912 bytes past with 933. The resolving window is a scan from the last anchor, the shape of
Corollary~\ref{cor:width}, met here on token sets with no maximum token length, where the corollary is silent. Inside
the conventional row's searched stretches the longest goes 104, 101, 91, 82, 44 and 16 bytes at widths $4$ to $128$.
Inside the block-comment row's it goes from 965 to 857 at 128 and to 127 at 1{,}024. Reach falls with the width and no
faster, so cutting a stretch down to the floor takes a window about as wide as the stretch. The vocabularies differ from
each other. The 384-merge vocabulary's 12 searched boundaries resolve by width 16, 11 of them by windows reaching no
anchor. The GPT-2 prefix's 6 resolve by width 16, every one by a window reaching back to its anchor, over 3 to 12 bytes.

Three neighbours frame the profile. Kaplan's pinch-points~\cite{kaplan2005tokenizing}, the concatenation check of
Lester, Ong and Sch\"afer~\cite{lester2016flow} and ZipLex's separability predicate~\cite{chassot2026ziplex} each
certify one join, of a text at a position or of two code strings, where a forced gap quantifies over every occurrence in
every completely tokenizable input. What buys reach past the floor is resolved context rather than boundaries: the
data-parallel simulation of Mytkowicz, Musuvathi and Schulte~\cite{mytkowicz2014dpfsm} and the bounded alternative
states of Barenghi et al.~\cite{barenghi2015parallel} enter a chunk in every live state and pay a scan per state kept.

\section{Evaluation}
\label{sec:evaluation}

The artifact runs the evaluation in the default test target and CI. The figures of this section are asserted rather than
merely printed, so a drifted number fails the test suite.

The random sweep is the output of \code{tools/probes/src/window\_gate.cpp} at munch commit
\code{707a79941472885a260c0bc96e615dd2fb5e99d2}, the commit at which the positive-width equivalent entered the library.
The v1.3.3 release the rest of this evaluation describes excludes the nullable sets and asserts the earlier counts. The
400 random token sets are on a three-symbol alphabet, generated by the probe itself with a pinned seed and draw order.
Of them, 266 are nullable and are decided through their positive-width equivalent, exactly as the artifact compiles
them, and 63 certify at least one byte exactly. Of the 337 that certify no byte, \textbf{326 gain a certified window
under the model, and 322 of those are witnessed}. For each witnessed grammar, the bounded search finds a completely
tokenizable input containing a certified window, with the covering token beginning at the reported origin. Each input is
verified as it is constructed, and the 322 aggregate is asserted. Another 4 model-positive grammars have no occurrence
within the bounded witness search and are reported as unresolved, never as rescues. The remaining 11 exhaust the
quotient with no window under the model, and none are inconclusive. Occurrence is a property of the concrete word rather
than its quotient key, so the witness search continues past the shortest certified length instead of stopping at the
first certifying word. Separate rewind-stress rows exercised 1{,}079{,}392 generated executions that scanned through the
window and contained at least one rewind, with zero disagreements against the shipped scanner. Of those executions,
418{,}466 tokenize their whole input completely and the remaining 660{,}926 have malformed suffixes past the window,
both counts asserted. The rewind-stress rows are an implementation stress check: the generated inputs were required to
scan through the window, not to tokenize completely. The random sweep additionally checked every certified two-byte
window over the probe's generated contexts. The length-one case reproduces the published certificate on all 400
grammars, the 266 nullable ones included.

Named token sets: all six exact-empty rows of the predecessor's applicability table, five C-like variants and JSON, gain
witnessed windows. One new cumulative C-like variant joins them as a seventh positive row. Table~\ref{tab:windows} shows
the concrete windows, with the retained key counts, one indicator of the search footprint, in place of the $6^{|Q^+|}$
bound. Beside the gate's six-status keys the table reports the keys of the three-status quotient and the nodes of the
chosen-origin search of Section~\ref{sec:smaller}, from the figure program \code{paper/figures/window\_search.cpp},
which runs both over the same token sets, asserts that their shortest lengths equal the gate's on every row, and pins
the counts. The example windows show the recovery anchors. The string and line-comment rows resynchronize at a newline
followed by a byte that must begin a token. The block-comment rows resynchronize at the byte pair \code{*/} followed by
whitespace, in comment context the closer, with the origin immediately after the pair. The certificate is
occurrence-universal, so it also covers occurrences where \code{*/} reads as two operator tokens. The pinned witnesses
exercise exactly that reading. The conventional row certifies at length two: its whitespace runs include the newline, so
\code{!} must begin a token there as well. The JSON row uses the RFC 8259~\cite{bray2017json} lexical forms over bytes
and assumes UTF-8 validity; it is not a conforming JSON processor. How often such windows occur in real corpora is an
empirical question for the measurement campaign, and no frequency claim is made here. The run $\code{a}^+$ is included
as the negative row, and $\code{a}^*$ would be decided as the same automaton. In that automaton every byte continues a
run as readily as it begins one. The model certifies no window, since absent-byte windows certify only vacuously under
Definition~\ref{def:window}, and the search exhausts its quotient, which is precisely the shape of a model-negative.

\begin{table}
\centering
\caption{Certified windows for the named token sets, from the gate's asserted rows: the shortest
model-certified length, one example window with its origin, and the quotient keys the search retained before
shortest-window stopping, under the six-status key of the gate and the three-status key of Section~\ref{sec:three}. The
last column is the nodes the chosen-origin search of Section~\ref{sec:chosen} holds when its shortest window appears, or
when the search exhausts. Every positive row's displayed window carries an asserted occurrence witness; among the
positive rows the cumulative one is new to this study, and every other positive row is an exact-empty row of the
predecessor's table.}
\label{tab:windows}
\small
\setlength{\tabcolsep}{3pt}
\begin{tabular}{lccrrr}
\toprule
 & & & \multicolumn{2}{c}{Retained keys} & \\
\cmidrule(lr){4-5}
Token set & \shortstack{Shortest\\model-\\certified $k$} & \shortstack{Example window\\at origin} & six-status &
three-status & \shortstack{Chosen-origin\\nodes} \\
\midrule
% The six-status column: tools/probes/src/window_gate.cpp; the two others: paper/figures/window_search.cpp.
C-like, string literals & 2 & \code{\char`\\n!} at 1 & 24 & 20 & 17 \\
C-like, line comments & 2 & \code{\char`\\n!} at 1 & 18 & 13 & 14 \\
C-like, block comments & 4 & \code{\char`\\t*/\char`\\t} at 3 & 53 & 30 & 30 \\
C-like, conventional & 2 & \code{\char`\\n!} at 1 & 27 & 22 & 16 \\
split-friendly, block comments & 4 & \code{\char`\\n*/\char`\\t} at 3 & 188 & 103 & 81 \\
C-like, cumulative (new here) & 4 & \code{\char`\\n*/\char`\\t} at 3 & 189 & 102 & 79 \\
JSON, RFC 8259 & 2 & \code{\char`\\t"} at 1 & 69 & 52 & 29 \\
$\code{a}^+$ (negative row) & none & search exhausted & 3 & 2 & 2 \\
\bottomrule
\end{tabular}
\end{table}

\section{Related work}
\label{sec:related}

The window generalizes the certified split point of~\cite{nidhogg2026splitpoints}, and inherits its relation to the
parallel lexing families. Composition carries every state and pays for it~\cite{mytkowicz2014dpfsm}. Speculation
predicts an entry state, validates, and re-executes on a miss~\cite{prabhu2010speculative}. Bounded alternative-state
scans disambiguate without single-state speculation~\cite{barenghi2015parallel}. Prescanning pays a pass over the
input~\cite{li2021plex}. The window, like the byte, is derived from the grammar before any input exists. A close
compiler-style neighbor in the maximal-munch setting is the streaming analysis of Li, Yang, and
Mamouras~\cite{li2026streaming}. The analysis statically computes a grammar's maximum token-neighbor distance and uses
the resulting bounded lookahead to emit a sequential maximal-munch stream without backtracking. The authors' scan
advances from a known token boundary, so the window decides maximality rather than recovering the origin of the token
covering an arbitrary occurrence, and parallelization is left there as future work. On the formal side of the same
setting, ZipLex formally verifies linear-time invertible maximal-munch lexing, with an abstraction capturing the
separability of tokens in a sequence~\cite{chassot2026ziplex}. Li and Mamouras formalize the uniform tokenization
problem and give $O(mn)$-time algorithms, linear in text length $n$ for a fixed grammar of size $m$. The algorithms
precompute what each suffix admits in a right-to-left pass before tokenizing from the input's
start~\cite{li2025uniform}. Neither line of work derives an occurrence-universal raw-window certificate or recovers the
covering token's origin from an arbitrary occurrence without scanning from a known boundary. Lester's boxing check
anticipates the flavor at a known join. Every lexer state possible after the first analyzed string must either already
emit a token, making the following character irrelevant to that string's lexing, or emit the same token immediately on
every character that may begin the second. Because of that requirement, the join is a lexeme boundary and concatenation
preserves token boundaries. The check is stated compactly in the position paper and in full in the journal
treatment~\cite{lester2013boxing, lester2016flow}. The check is per-join over two analyzed fragments rather than a
certificate over every occurrence of a grammar-derived word, and it does not recover a covering origin. Recent
LLM-tokenizer work addresses input-specific seams instead. LoPT validates position-aligned tokenizations of overlapping
chunks and adjusts chunk length where needed~\cite{shao2026lopt}. For BPE, recent work bounds the streaming delay of
ordered merge rules from a known beginning~\cite{mamouras2026bpedelay} or maintains the tokenization incrementally over
every prefix~\cite{jiang2026incremental}. Hayase, Liu, Smith, and Oh enumerate, for a byte prefix of BPE-tokenized text,
the tokenizations whose last token straddles the prefix end, a valid covering tree over one concrete
input~\cite{hayase2025bytesampler}. The certificate here is grammar-derived and occurrence-universal rather than
input-relative and enumerative, and it fixes one origin for every occurrence. The parsing side of that pipeline is
active again. Cyclic operator precedence grammars~\cite{chiari2025cyclic} admit parallel-suitable chunking of flat
unbounded terminal-string substructures, with precedence relations defined between adjacent terminals. In a conventional
pipeline those terminals are lexer output, so that application presupposes tokenization, though the formalism itself is
not restricted to pre-tokenized input. A classical antecedent of window-determines-state is the definite automaton,
studied in depth by Perles, Rabin, and Shamir~\cite{perles1963definite}. The definite automaton's operational form for a
fixed $k$ is $k$-locality: any $k$ consecutive symbols force a unique state, every word of length $k$ being
synchronizing~\cite{holub2009parallel}. Both quantify uniformly over all windows and speak of states. The certificate
here is per-window, speaks of token boundaries under maximal munch, and recovers an origin. Raw-state synchronization
alone does not identify the start of the covering maximal-munch token. The classical special case of boundary recovery
is code synchronization. Proposition~\ref{prop:codes} states the relation exactly for prefix codes, where left-to-right
tokenization realizes the code factorization. Over a prefix code the certified windows that occur are the synchronizing
splits whose right half sits inside one codeword. A synchronizing pair~\cite{berstel2010codes} with a nonempty right
component yields, on its concatenation, a certified window whose origin is the start of the right component's last
codeword. An early member of that family is the comma-free code of Golomb, Gordon and Welch~\cite{golomb1958commafree}.
A comma-free code is a block code no codeword of which occurs across the boundary of two adjacent codewords, so every
codeword is a certified window with origin $0$ for the code's own factorization. Finite synchronization
delay~\cite{restivo1975synchronization} bounds how many codewords a synchronizing pair needs, hence a byte bound for a
finite code. The $ww$-style resumption argument survives into the partial-DFA treatment~\cite{berlinkov2021partial}. The
explicit modern bounded-window form of that special case is the synchronizing morphism. There a window of bounded length
suffices to detect boundaries between codewords~\cite{fici2025morphisms}, in a morphic code-factorization setting rather
than among competing prioritized token languages under maximal munch. Uniquely decipherable codes may share prefixes.
What such codes guarantee is a unique factorization, with no maximal-munch priority resolving overlaps between competing
token languages. The difference is where the origin machinery here earns its existence. The relationship to reset words
is one-way and stops at length one. A useful certified byte induces a reset-like action on the partial live automaton
with domain $\{q_0\}$, under the stated re-entrancy qualification. The correspondence fails under the classical
complete-DFA reading~\cite{volkov2008synchronizing}. The correspondence does not extend. A certified window need not be
a reset word of the token DFA at all, since over $\{\code{a}, \code{b}\}$ the window \code{ab} certifies at origin 1
while the action of \code{ab} on the partial automaton is empty. A rank-one letter need not certify either, since over
the token language $\code{b}^*\code{a}$ the letter \code{b} can act with rank one while $q_0$ is re-entrant and \code{b}
occurs inside tokens. Certified windows synchronize token-origin information in an enriched cloud model; they need not
synchronize the raw token DFA. The complexity of the neighborhood is known. Checking careful synchronizability of a
partial automaton, and finding a shortest carefully synchronizing word, are PSPACE-complete already over two-letter
alphabets~\cite{martyugin2010careful}. Here careful means the word stays defined from every state and maps all states to
one. Section~\ref{sec:complexity} proves the model's existence problem PSPACE-complete by a different route, from the
intersection non-emptiness of DFAs~\cite{kozen1977lower}. The two problems differ in what a word must do: a careful
reset word maps every state of the automaton to one, where a certified window need not synchronize the token DFA's
states, only make the cloud's survivors agree on an origin. The per-grammar retained-key counts, one footprint indicator
rather than a cost model, stayed far below the worst case throughout.

Two practices are the practical counterparts. Compiler panic-mode recovery discards input to a recovery
set~\cite{aho2006compilers}, which may itself be grammar-derived. The distinction is post-error parser recovery against
a pre-input lexical boundary guarantee. Incremental lexing in the style of Wagner and Graham restarts from per-token
scanner-state snapshots in a versioned token stream, with dynamically maintained lookahead dependencies, exact for the
text they were cached against~\cite{wagner1997lexing}. A certified window is grammar-universal instead, valid in every
completely tokenizable context and known before any input exists, at the price of existing only where the token set
admits one.

Very recent concurrent work establishes nearby but different local guarantees. TokTier certifies input-relative splice
junctions for selected frozen pre-tokenizer and BPE pipelines. TokTier matches cached against fresh runs, uses
family-specific synchronizing character-class transitions proved to reset the pre-tokenizer under every left context,
and proves exact recombination of its overlapping windows~\cite{zhang2026toktier}. ReTokSync monitors the receiver-view
tokenization of one concrete generated stream and triggers a corrective reset when ambiguity
occurs~\cite{wang2026retoksync}. Borsotti, Crespi Reghizzi, and Pradella define precedence relations over synthesized
attributes for a deterministic parser using one-symbol left and right neighborhoods~\cite{borsotti2026attribute}. None
decides occurrence-universal bounded words for competing prioritized token languages under maximal munch, nor recovers
the origin of the token covering an arbitrary cut.

Two formal treatments of tokenization itself address a different question. Kaplan models a language's tokenizing
conventions as a finite-state transduction and emits output at input-relative pinch-points. Pinch-points are positions
where all live analysis paths converge to a single transducer state, so settled prefixes stream out as a particular text
is processed~\cite{kaplan2005tokenizing}. Cognetta and Okazaki encode the tokenizations of a regular language as
finite-state transduction and represent canonical MaxMatch and BPE subword tokenizers as
transducers~\cite{cognetta2025fst}. Neither derives occurrence-universal windows from the grammar alone, and neither
recovers the origin of the covering maximal-munch token at an arbitrary occurrence.

Symbolic dynamics comes closest in shape. A resolving block is a block of a factor subshift all of whose admissible
preimages agree at a selected coordinate~\cite{adler1983sliding, marcus1985sofic}. Such a block is an
occurrence-universal lift of a local observation to hidden state, introduced as a means of resetting an encoding
automaton in sliding-block code construction. A certified window shares that shape, reading the window as the observable
block and the covering token's start as the hidden coordinate. The classical theory, however, is stated for factor maps
of subshifts, without token priorities, maximal munch, backup, or any tokenization semantics, and it neither defines nor
decides the lexical instantiation.

The areas surveyed above are definite and local automata, synchronizing words and codes under both the complete and the
partial reading, resolving blocks, careful synchronization, streaming, verified, and uniform maximal-munch tokenization,
incremental relexing, the parallel lexing families, and the parallel parsing line. In the prefix-code and
code-factorization special case, synchronizing pairs and bounded windows already recover codeword
boundaries~\cite{restivo1975synchronization, fici2025morphisms}. Beyond that case, we have not found prior work in these
areas that defines or decides in general, for competing prioritized token languages under maximal munch, a
grammar-derived bounded word whose every occurrence identifies the origin of the covering token at an arbitrary cut. Nor
have we found, in these areas, prior work instantiating resolving-block or local-decoding machinery for prioritized
maximal-munch token origins, nor a derivation of this conservative compiled-table cloud model. The closest results
decide maximality from a known boundary, bound retained memory, assume the boundaries they schedule, recover a state
without an origin, or validate a concrete overlap or splice after local retokenization. We are likewise not aware of an
implementation that derives certified windows from a compiled token set and checks them against a running scanner, other
than the instruments this paper reports and their supported successor in release 1.4.0.

\section{Limitations}
\label{sec:limitations}

All deliberate. The soundness proof speaks only of completely tokenizable inputs: malformed input is outside the proof,
and no consumed-prefix analogue is claimed. Negatives are model-relative: the model refuses windows a greedy scanner
would allow, so an exhausted search means no window \emph{under this model}, never that none exists. The worst case is
exponential, and Section~\ref{sec:complexity} shows that no polynomial algorithm for the model's existence question
exists unless $\mathrm{P} = \mathrm{PSPACE}$, with shortest windows themselves of length $2^{\Omega(\sqrt n)}$ on a
family of token sets. The probe and the shipped search are budgeted, though the retained keys stayed below 200 on every
named row, at most 32 with mean 9.9 among the 337 no-byte grammars. The complexity results concern the model and the
compiled automaton; the complexity of Definition~\ref{def:window} itself, and bounds in the size of a regular
expression, are not claimed. The decision of global completeness has an EXPSPACE upper bound and no matching lower
bound, and the exact witnessed synthesis was not run over the random sweep. The evaluated v1.3.3 artifact ships no
window-planning API, and its probe excludes the nullable sets, so the sweep of Section~\ref{sec:evaluation} is
reproduced at the later commit named there. The artifact's certificate machinery is a probe, on the principle that the
certificate's formulation should freeze in this paper before becoming a public contract. Release 1.4.0 made the
formulation a public contract, \code{is\_split\_window()}. The sweep at the later commit asks that contract beside its
own model, at the coverage the sweep's source states. The coverage is every one-byte answer of the named rows and of the
400 random token sets, every model-positive random two-byte answer, and the named certificates and refusals, with the
count of those checks pinned. The window planner the same release added beside the decision lies outside this paper's
evaluation. The real-corpus evidence is the exploratory campaign of Section~\ref{sec:profile}, recorded outside the
artifact and asserted by no test. Window occurrence frequency on representative corpora remains a question for a
controlled campaign.

\section{Conclusion}
\label{sec:conclusion}

A certified split window $(W, o)$ of a compiled token set is a byte string such that in every completely tokenizable
input containing $W$, the token covering the occurrence's final byte begins exactly $o$ bytes into it. Model
certification is decided from the compiled tables alone by running a conservative cloud of token-prefix hypotheses
across the window, seeding new trajectories only where the automaton accepts, and demanding unanimity on an in-window
origin. The soundness argument is a representation invariant: the final segmentation's actual token-prefix history is
always among the hypotheses, so unanimity can only land on the truth. Maximal-munch backup never needs simulating,
because the model tracks where tokens begin rather than where the read head wanders.

The search is a decision procedure for the model, not a heuristic. A quotient of the reachable clouds, exact under the
stated single-occupancy invariant, makes breadth-first search terminate by exhaustion, so every answer is a certified
window with its origin or a proof that the model admits none at any length. Two smaller searches decide the same
question, a three-status quotient and an unambiguous automaton that follows one chosen origin, the latter shipped as
\code{shortest\_split\_window()} by release 2.2 of munch. The exponential is not the quotient's fault: existence of a
model certificate is PSPACE-complete, plain or witnessed, and a shortest window can be $2^{\Omega(\sqrt n)}$ bytes
long over $n$ states. The cloud's hypotheses are exactly the histories of arbitrary factorizations of words of the
token language. That is why the vacuous window is certified, why token kinds do not matter to the search, why the
model is exact on the polynomially decidable class where every factorization is the scan's, and how a shortest
witnessed window and the model's global completeness are decided. The three-layer honesty is deliberate and
permanent. Model-positive answers are semantically certified within the stated flat, completely-tokenizable scope.
Negatives are model-relative: the model refuses windows a greedy scanner would allow, and the two asserted witnesses
exhibit the over-approximation. The strictness corollary locates the over-approximation's minimum nonvacuous length at
two, since at length one the model provably coincides with the published predicate, which the predecessor's necessity
theorem makes exact for occurring bytes. Over a prefix code the certificate is code synchronization under another name,
the occurring certified windows being the synchronizing splits whose right half sits inside one codeword. The named
token sets studied here are not codes. A token set in which some token matches the empty string is decided through its
positive-width equivalent, which unrolls the accepting start state and changes no scan, so nothing is set aside on that
account.

The generalization does what it was built for: it recovers every exact-empty row of the predecessor's table, all
witnessed. Of the 337 random token sets certifying no byte, 322 gain a witnessed certified window with zero inconclusive
searches. The predecessor's six exact-empty rows gain witnessed windows of two to four bytes, with the concrete windows
in Table~\ref{tab:windows}. Separate rewind-stress rows exercised 1{,}079{,}392 generated executions that scanned
through the window and contained at least one rewind, with zero disagreements against the shipped scanner. The figures
are printed and asserted in the artifact, so they fail the test suite if they move. What this paper deliberately does
not deliver is the cut itself. The v1.3.3 artifact evaluated here ships a planner over single-byte certificates only,
itself not evaluated here. The window planner that release 1.4.0 added beside the decision is outside this evaluation.
Turning a window's boundary into a parallel cut is an explicit composition with the predecessor's prefix-stability
result. Whether windows occur often enough in real corpora to plan balanced chunks is an empirical question on which
this paper relies on no controlled evaluation. The artifact archives exploratory preview measurements only, and the
anchor densities of Section~\ref{sec:profile} are an exploratory campaign outside the artifact. A controlled evaluation,
on representative corpora with planning times, occurrence frequencies, and end-to-end comparisons against byte
certificates, is the natural next step and is not claimed here.

Read gap by gap, the certificate is one pattern of a window's boundary profile: one forced gap followed by crossed gaps
to the window's end. A rule that cuts at a window's gap and reconciles nothing is sound exactly at forced gaps, so the
profile is the whole anchor supply a window of its width can give such a rule. Over a token set with a longest token, a
forced gap propagates: the scan from it forces every boundary it commits inside the window, so a boundary resolves at a
width about its distance to the last anchor. No single width decides every gap of the JSON row, or of the C-like rows
with strings or comments, since no bounded window sees whether an unmatched quote or an open comment lies to its left.
At the campaign width the profile adds anchors on JSON and the byte-pair vocabularies, none on the C-like rows, and
shortens no longest stretch, whose floor is the longest token.

The contribution is therefore narrow and exactly bounded, in the same sense as its predecessor: not another way to
recover context, but the certification step, extended from single bytes to bounded windows, with the origin recovered
and the conservatism exhibited. Certified windows synchronize token-origin information in an enriched cloud model, and
need not synchronize the token DFA itself. The predecessor's one-way reset-word connection applies at length one, and no
general implication extends to longer windows.

\section*{Tools}

Large language models assisted with drafting, with checking citations against primary records and with reviewing the
implementation. The author verified each suggestion by derivation, against primary sources, or against the
implementation, its tests and the archived artifacts, and is responsible for all content.

\bibliographystyle{plainnat}
\bibliography{refs}

\end{document}
